\section{Background}
\label{sec:background}

In this section, we provide some necessary mathematical preliminaries and establish nomenclature and notation. Further background is provided in Appendix~\ref{app:further-background}, which readers unfamiliar with spinors, filtrations, $\ZZ$-gradings on superalgebras, Lie superalgebra cohomology or homogeneous spaces, or those who require a review, may wish to read first.

\subsection{Flat model algebras}
\label{sec:flat-model-algs}

The Euclidean and Poincar\'e algebras are the isometry algebras of Euclidean and Minkowski space respectively. They can be considered as ``flat models" for isometry algebras in Riemannian and Lorentzian geometry in the informal sense that they are the isometry algebras of the maximally symmetric flat spaces in their respective settings. A more precise meaning can be given to this using Cartan geometry, in which a geometric structure on a manifold is viewed as a local equivalence to some homogeneous space -- see \cite{Wise2010} for a pedagogical introduction emphasising links to pseudo-Riemannian geometry and application to general relativity in particular and \cite{Cap2009_1} for a more systematic treatment. Work by \v{C}ap and Neusser \cite{Cap2009} shows that, under certain regularity conditions, one can assign a Lie algebra of symmetries to a Cartan geometry, and that this algebra is a \emph{filtered deformation}\footnote{Filtrations are discussed in Appendix~\ref{sec:filtered-alg} and filtered deformations in Section~\ref{sec:filtered-def}. The term ``filtered deformation'' is not used in \cite{Cap2009} but was previously introduced by Cheng and Kac \cite{Cheng1998}.}
of a subalgebra of the flat model algebra. Pseudo-Riemannian manifolds with signature $(p,q)$ can be considered as Cartan geometries for the local model $\ISO(p,q)/\Orth(p,q)$ with a Cartan connection encoding both the spin connection and a soldering form. If this spin connection is Levi-Civita, the symmetry algebra of the Cartan connection is precisely the Lie algebra of infinitesimal isometries (Killing vectors), and the result discussed above says that it is filtered, with its associated graded Lie algebra being a subalgebra of $\fiso(p,q)$.

Although we do not use the formalism of Cartan geometry in this work, we can view the result on the structure of Killing superalgebras already mentioned in the introduction as a ``super'' generalisation of this phenomenon in pseudo-Riemannian geometry, and it would be subsumed by a generalisation of the \v{C}ap--Neusser result \cite{Cap2009} to super-Cartan geometry. Our (much more explicit) version of this result is Theorem~\ref{thm:killing-algebra-filtered}, which tells us that Killing (super)algebras are filtered subdeformations of the $\ZZ$-graded (super)algebras we will introduce in this section. Motivated by the connection to Cartan geometry, we will refer to these $\ZZ$-graded (super)algebras as \emph{flat model $($super$)$algebras}.

See Appendix~\ref{sec:inner-products-spinors} for further background on inner product spaces and related structures including Clifford algebras and spinors.

\subsubsection{Isometry algebra of an inner product space}

Let $(V,\eta)$ be a (pseudo-)inner product space. We denote by $\fiso(V,\eta)$ the isometry algebra of (the affine space modelled on) $(V,\eta)$ and by $\fso(V,\eta)$ the Lie algebra of $\eta$-skew-symmetric endomorphisms of $V$, although we usually suppress $\eta$ in the notation. As a vector space $\fiso(V)=V\oplus\fso(V)$, and it is equipped with the Lie bracket
\begin{align}\label{eq:flat-bracket}
 \comm{A}{B} = AB - BA,\qquad
 \comm{A}{v} = Av,\qquad
\comm{v}{w} = 0,
\end{align}
where $A,B\in\fso(V)$ and $v,w\in V$. In the Euclidean case $\fiso(V)$ is known as the \emph{Euclidean algebra}; in the Lorentzian case it is the \emph{Poincar\'e algebra} and $\fso(V)$ is the \emph{Lorentz algebra}.

We write $\fiso(p,q):=\fiso(\RR^{p,q})$ where $\RR^{p,q}$ is the standard inner product space with signature~$(p,q)$; if $\eta$ has signature $(p,q)$, the isomorphism $(V,\eta)\cong\RR^{p,q}$ induces a Lie algebra isomorphism $\fiso(V)\cong\fiso(p,q)$.

\subsubsection{Isometry algebra with spinors}
\label{sec:isom-alg-spinors}

We now construct a $\ZZ$-graded Lie algebra or Lie superalgebra\footnote{See Appendix~\ref{sec:graded-superalg} for a definition.} $\fs=\fs_{-2}\oplus\fs_{-1}\oplus\fs_{0}$ whose even part is $\fiso(V)$ as follows. Let
$\fs_0 = \fso(V)$, $
 \fs_{-1} = S$, $
 \fs_{-2} = V$, $
 \fs_i = 0$ otherwise,
where $(V,\eta)$ has signature $(p,q)$ and $S$ is a (possibly $N$-extended) spinor module; that is, we take
\begin{itemize}\itemsep=0pt
	\item $S = \ssS_N := N\ssS_1$ for some integer $N\geq 1$ for $p-q\neq 0,4\mod 8$, where $S_1$ is the unique irreducible real spinor module of $\Cl_{\overline 0}(p,q)$,
	\item $S = \ssS_{(N_+,N_-)} = N_+\ssS_+\oplus N_-\ssS_-$ for integers $N_+, N_-\geq 0$ not both zero if $p-q= 0,4\mod 8$ where $S_\pm$ are the two (non-isomorphic) irreducible real spinor modules of $\Cl_{\overline 0}(p,q)$~-- without loss of generality, we can take $N_+\geq N_-$.
\end{itemize}
We have used conventions for spinor modules explained in Appendix~\ref{sec:pinor-spinor}. We let the bracket on $\fs_{\overline 0}=\fiso(V)$ be the one defined by \eqref{eq:flat-bracket} and note that this respects the $\ZZ$-grading. In order to extend this to a Lie (super)bracket $\fs\otimes\fs\to\fs$ in a way which respects the $\ZZ$-grading, we must specify maps $\fs_{-1}\otimes\fs_0\to\fs_{-1}$ and $\fs_{-1}\otimes\fs_{-1}\to\fs_{-2}$ -- all other components must be zero. We take the first of these maps to be the natural action of $\fso(V)$ on $S$, so the brackets are given by
\begin{align}\label{eq:flat-spinor-bracket}
 \comm{\epsilon}{\zeta} = \kappa(\epsilon,\zeta),\qquad
 \comm{A}{\epsilon} = A\cdot\epsilon,\qquad
 \comm{v}{\epsilon} = 0,
\end{align}
for $A\in\fso(V)$, $v\in V$ and $\epsilon,\zeta\in S$, where we leave the map $\kappa\colon \Wedge^2 S\to V$ (for a Lie algebra) or~${\kappa\colon \Odot^2 S\to V}$ (for a Lie superalgebra) to be determined. The action of $A\in\fso(V)$ on $\epsilon\in S$ is defined in Appendix~\ref{sec:pin-spin-group} (see equation~\eqref{eq:sov-spinor-action}).

It remains to check the Jacobi identities involving elements of $\fs_{-1}$. The $[\fs_0,\fs_0,\fs_{-1}]$ identity is satisfied precisely because the $\fs_0\otimes\fs_{-1}\to\fs_{-1}$ component of the bracket is a module action. The $[\fs_0,\fs_{-1},\fs_{-2}]$ and $[\fs_{-2},\fs_{-1},\fs_{-1}]$ identities are trivial since $\comm{\fs_{-1}}{\fs_{-2}}=\comm{\fs_{-2}}{\fs_{-2}}=0$. The remaining component is $[\fs_0,\fs_{-1},\fs_{-1}]$, which is satisfied if and only if $\kappa$ is $\fso(V)$-equivariant. Thus we make the following definition.

\begin{Definition}[flat model (super)algebra, Dirac current]\label{def:flat-model-alg}
	Let $(V,\eta)$ be an inner product space and $S$ an (extended) real spinor module of $\fso(V)$. Then
	\begin{itemize}\itemsep=0pt
		\item A \emph{flat model algebra} is a $\ZZ$-graded Lie algebra $\fs$ with underlying $\ZZ$-graded vector space $V\oplus S\oplus \fso(V)$ and bracket defined by \eqref{eq:flat-bracket} and \eqref{eq:flat-spinor-bracket} for some $\fso(V)$-equivariant map $\kappa\colon \Wedge^2 S\to V$.
		\item A \emph{flat model superalgebra} is a $\ZZ$-graded Lie superalgebra $\fs$ with underlying $\ZZ$-graded vector space $V\oplus S\oplus \fso(V)$ and bracket defined by \eqref{eq:flat-bracket} and \eqref{eq:flat-spinor-bracket} for some $\fso(V)$-equivariant map $\kappa\colon \Odot^2 S\to V$.
	\end{itemize}
	In either case, $\kappa$ is referred to as the \emph{Dirac current}.
	
	Furthermore, we say that $\fs$ is \emph{minimal} if $S=\ssS_1=\ssS$ or $\ssS_{(1,0)}=\ssS_+$ and that it is \emph{$N$-extended} if $S=\ssS_N$ for $N>1$ or $S=\ssS_{(N_+,N_-)}$ for $N=N_++N_->1$; in the latter case we also say it is \emph{$(N_+,N_-)$-extended}.
\end{Definition}

We note that all flat model (super)algebras are referred to as \emph{$N$-$($super$)$-extended Poincar\'e algebras}, \emph{$\pm N$-extended Poincar\'e algebras} \cite{Alekseevsky1997} or \emph{Poincar\'e Lie superalgebras} \cite{Cortes2020,Gall2021}, but we reserve the name \emph{Poincar\'e} for Lorentzian signature. Our $\ZZ$-grading convention follows that used in previous work on Spencer cohomology and supergravity \cite{Beckett2021,deMedeiros2016,deMedeiros2018,Figueroa-OFarrill2017_1,Figueroa-OFarrill2017}.

Since inner product spaces of equal signature are isomorphic, the data required to specify a~flat model (super)algebra are the signature $(p,q)$, the integer $N$ or pair of integers $(N_+,N_-)$ and the Dirac current $\kappa$. The possible Dirac currents were classified by Alekseevsky and Cort\'es \cite{Alekseevsky1997} by reducing the problem to a classification of $\fso(V)$-invariant bilinears on the irreducible pinor module $\ssS$. More recently, this classification has been improved (at least in the superalgebra case) by describing a necessary and sufficient condition for two different Dirac currents to define isomorphic superalgebras \cite{Cortes2020}, and the possible $R$-symmetry groups -- a particular group of Lie superalgebra automorphisms which serves as an algebraic invariant -- for a wide class of Dirac currents have been described, the comparison of which provides a sufficient condition to distinguish non-isomorphic algebras which is simpler to check \cite{Gall2021}.

Our main focus in this work, or at least its latter half, will be on Lorentzian signature and ($N$-extended) Poincar\'e \emph{super}algebras, mainly due to important technical simplifications which arise in this case from the homogeneity theorem (see Theorem~\ref{thm:homogeneity}). Nonetheless, we will provide a significant amount of formalism in general signature. We will not consider conformal superalgebras or central extensions, nor any other generalisation, although flat model algebras extended by $R$-symmetry will be considered in upcoming work. There is a significant amount of literature on all of these generalisations: for the Lorentzian case, see the classic work of Strathdee (which also discusses the representation theory of the superalgebras) \cite{Strathdee1986} and the review of Van Proeyen \cite{VanProeyen1999}; for general signature, see \cite{Alekseevsky1997,Alekseevsky2005,DAuria2001_1} and also the more recent work \cite{Gall2021}.

\begin{Remark}
	The Dirac current can often be seen as a component of a polyvector-valued (or, dually, polyform-valued) \emph{squaring map} which has been used to study generalisations of Killing spinors algebro-geometrically \cite{Cortes2021}. The other components of this map are often called \emph{Dirac bilinears} or \emph{Dirac $p$-forms}; in the $N$-extended case, these forms typically ``carry internal indices'', i.e., they are valued not in $\Wedge^p V$ for some $p$ but in $\Wedge^p V\otimes\Delta$ where $\Delta$ is some auxiliary space.
	
	The extension of the (super)algebras described above to include these higher forms in the odd-odd components of the bracket has been considered, e.g., in~\cite{Alekseevsky2005,Figueroa-OFarrill2009}. The so-called FDAs (``free'' differential (graded) algebras) of the rheonomic formalism of supergravity \cite{DAuria1982,DAuria1980}, which have been put on a rigorous mathematical footing as $L_\infty$-algebras \cite{Schreiber2015,Schreiber2009,Stasheff2018}, also incorporate these forms, albeit in a different manner.
	All of these extensions can provide information about higher gauge theory data on supergravity backgrounds which is unfortunately invisible to the Lie (super)algebras discussed in this work. We hope to understand how the perspective proposed here can be extended to include such gauge-theoretic information in future work -- and we note that the natural setting for doing so would be higher Cartan geometry as in \cite{Schreiber2009} -- but such considerations are beyond our current scope.
	
	Notwithstanding the comments above, note that although higher Dirac forms are not used in our flat model (super)algebras, they nonetheless do appear in explicit descriptions of their filtered deformations (such as those found in \cite{Beckett2019,deMedeiros2016,deMedeiros2018,Figueroa-OFarrill2017} and the companion work \cite{Beckett2024_3}), as they provide a convenient way to parametrise bilinear maps on the spinor module (e.g., the maps denoted $\gamma$ appearing in \eqref{eq:22-cocycle-comps} and \eqref{eq:def-brackets-general}).
\end{Remark}

\subsection{Filtered deformations and Spencer cohomology}
\label{sec:spencer-thy}

Here we present some background material on the Spencer cohomology and filtered deformations of $\ZZ$-graded Lie superalgebras. Our main source for this section is the work of Cheng and Kac~\cite{Cheng1998}. The base field is assumed to be either $\RR$ or $\mathbb{C}$. We recommend that readers unfamiliar with gradings and filtrations of vector (super)spaces and (super)algebras read Appendix~\ref{sec:grading-filtration-superspace} before this section.

\subsubsection{Spencer cohomology}
\label{subsec:spencer-cohomology}

Let $\fg=\bigoplus_{k=-h}^\infty\fg_k$ be a $\ZZ$-graded Lie superalgebra of finite depth\footnote{The Spencer cohomology theory in its original form was only applicable to algebras of depth 1; for those of greater depth Cheng and Kac use the term \emph{generalised} Spencer cohomology. We do not make this distinction.}
$h$ such that $\dim\fg_k<\infty$ for all $k$. Ultimately, we will be interested in the case where $\dim\fg<\infty$, but our source \cite{Cheng1998} mainly treats examples where $\fg$ itself is infinite-dimensional, and the theory presented here applies in any case. We let $\fg_-=\bigoplus_{k<0}\fg_k$ and note that it is a finite-dimensional subalgebra.

%%%%%%%%%% {The Spencer complex}

We define the Spencer cochain complex $\qty(\ssC^\bullet(\fg_-;\fg),\partial^\bullet)$ as follows. For $p<0$, we set $\ssC^p(\fg_-;\fg)\allowbreak:=0$, while for $p\geq 0$, we have
$\ssC^p(\fg_-;\fg) := \qty(\Wedge^p \fg_-^*) \otimes \fg \cong \Hom\qty(\Wedge^p \fg_-,\fg)$,
where $\Wedge^p$ is meant in the super-sense (see equation \eqref{eq:super-wedge}). The Spencer differential $\partial^\bullet\colon \ssC^\bullet(\fg_-;\fg)\to \ssC^{\bullet+1}(\fg_-;\fg)$ is given by the following formula: for $\phi\in \ssC^p\qty(\fg_-;\fg)$, $X_1,\dots, X_r\in\fg_{\overline{0}}$ and $Y_1,\dots,Y_s\in\fg_{\overline{1}}$, where $r+s=p+1$,
\begin{gather}
	(\partial\phi)\qty(X_1,\dots,X_r,Y_1,\dots,Y_s)\nonumber\\
\qquad	= \sum_{i=1}^r (-1)^{i+1} [X_i,\phi\bigl(X_1,\dots,\hat{X}_i,\dots,X_r,Y_1,\dots,Y_s\bigr)]	\nonumber\\
\phantom{\qquad=}{}+ (-1)^{r} \sum_{s=1}^r [Y_j,\phi \bigl(X_1,\dots,X_r,Y_1,\dots,\hat{Y}_j,\dots,Y_s)]\nonumber	\\
\phantom{\qquad=}{}+ \sum_{1\leq i<j\leq r} (-1)^{i+j} \phi \bigl(\comm{X_i}{X_j},X_1,\dots,\hat{X}_i,\dots,\hat{X}_j,\dots,X_r,Y_1,\dots,Y_s\bigr)\nonumber	\\
\phantom{\qquad=}{}+ \sum_{i=1}^r\sum_{j=1}^s (-1)^{i} \phi \bigl(X_1,\dots,\hat{X}_i,\dots,X_r,\comm{X_i}{Y_j},Y_1,\dots,\hat{Y}_j,\dots,Y_s\bigr)	\nonumber\\
\phantom{\qquad=}{}+ \sum_{1\leq i<j\leq s} \phi \bigl(\comm{Y_i}{Y_j},X_1,\dots,X_r,Y_1,\dots,\hat{Y}_i,\dots,\hat{Y}_j,\dots,Y_s\bigr),\label{eq:CEdiffLiesuper}
\end{gather}
where the hat decorations denote omission of an entry. It can be checked that $\partial^2=0$, and so $\qty(\ssC^\bullet(\fg_-;\fg),\partial^\bullet)$ is indeed a cochain complex. This is nothing but the standard Chevalley--Eilenberg complex for $\fg_-$ with values in $\fg$, where the former acts on the latter via the restriction of the adjoint representation. However, since $\fg$ is graded, there is additional structure.

The grading on $\fg$ induces a grading on $\fg^*$ defined by $(\fg^*)_k:=(\fg_{-k})^*$. This in turn induces a grading on the space of $p$-cochains, in the usual manner for tensor products of graded spaces. Equivalently, we can consider a $p$-cochain to be homogeneous of degree $d$ if it has degree $d$ as a~map $\Wedge^p \fg_-\to\fg$. We then denote the space of $p$-cochains of degree $d$ as $\ssC^{d,p}(\fg_-,\fg)$,
\begin{equation}\label{eq:spencer-complex-grading}
	\ssC^{d,p}(\fg_-;\fg) = \qty(\Wedge^p \fg_-^*\otimes \fg)_d = \Hom\qty(\Wedge^p \fg_-;\fg)_d,
\end{equation}
so that
\[
	\ssC^p(\fg_-;\fg) = \bigoplus_{d\in\ZZ} \ssC^{d,p}(\fg_-;\fg) \qquad\text{and} \qquad\ssC^{d,\bullet}(\fg_-;\fg) := \bigoplus_{p\in\ZZ} \ssC^{d,p}(\fg_-;\fg).
\]
This grading is preserved by the Spencer differential, so by restricting the differential to the subspaces of cochains with degree $d$, for each $d\in\ZZ$ we have a subcomplex $\bigl(\ssC^{d,\bullet}(\fg_-;\fg),\partial^\bullet\bigr)$.

Finally, we note that the adjoint action of $\fg_0$ on $\fg$ preserves the grading, thus so does the natural action of $\fg_0$ on cochains induced by the adjoint representation (which also trivially preserves the homological grading); each $\ssC^p(\fg_-;\fg)$ is a graded $\fg_0$-module. This action also preserves the differential and thus all of the structure on the subcomplex $\bigl(\ssC^{d,\bullet}(\fg_-;\fg),\partial^\bullet\bigr)$.

%%%%%%%%%% {The cohomology}

We define the cocycles, coboundaries and cohomology of the $d$-subcomplex in the usual way
\begin{gather*}
	\ssZ^{d,p}(\fg_-;\fg) = \ker\qty(\partial^p\colon \ssC^p(\fg_-;\fg)\to \ssC^{p+1}(\fg_-;\fg)),	\\
	\ssB^{d,p}(\fg_-;\fg) = \Im\qty(\partial^p\colon \ssC^{p-1}(\fg_-;\fg)\to \ssC^{p}(\fg_-;\fg)),	\qquad
	\ssH^{d,p}(\fg_-;\fg) = \faktor{\ssZ^{d,p}(\fg_-;\fg)}{\ssB^{d,p}(\fg_-;\fg)},
\end{gather*}
and these are all naturally $\fg_0$-modules, since the action of $\fg_0$ preserves $\bigl(\ssC^{d,\bullet}(\fg_-;\fg),\partial^\bullet\bigr)$. Of course, the space of $p$-cocycles of the whole complex is simply the direct sum of the graded $(d,p)$-cocycles, and similarly for the coboundaries, whence we have a natural isomorphism of $\fg_0$-modules
\[
	\ssH^p(\fg_-;\fg) \cong \bigoplus_{d\in\ZZ}\ssH^{d,p}(\fg_-;\fg),
\]
where the former is the cohomology of the whole complex.

The definition of the differential \eqref{eq:CEdiffLiesuper} on its own is not very enlightening, so let us demonstrate it explicitly on cochains of low homological degree and give and interpretation of the spaces of cocycles and the cohomologies. Fix $d\in\ZZ$ and homogeneous elements $X,Y,Z\in\fg_-$; we denote their degree by $|X|$ etc. Then for $\phi\in\ssC^{d,0}(\fg_-;\fg)=\fg_d$, we have
$(\partial \phi)(X) = \comm{X}{\phi}$,
which shows that $(d,0)$-cocycles are $\fg_-$-invariants in $\fg_d$; equivalently, they are degree-$d$ elements of $\fg$ which centralise $\fg_-$, and $(d,1)$-coboundaries are (super)derivations $\fg_-\to\fg$ given by restricting inner derivations of $\fg$ induced by elements of degree $d$. For $\phi\in\ssC^{d,1}(\fg_-;\fg)=\Hom(\fg_-;\fg)_d$,
\[
	(\partial \phi)(X,Y) = \comm{X}{\phi(Y)} - \comm{Y}{\phi(X)} -\phi(\comm{X}{Y}),
\]
so $(d,1)$-cocycles are (super)derivations $\fg_-\to\fg$ of degree $d$. Thus we find that $\ssH^{d,0}(\fg_-;\fg) \cong (\fg_d)^{\fg_-}$ and $\ssH^{d.1}(\fg_-;\fg)=\operatorname{der}(\fg_-,\fg)_d/\fg_-$. For cocycles in higher degree, we do not have such simple descriptions, but for the calculations in this work we will need expressions only for homological degree $p\leq 2$. Thus we make the $p=2$ case our final example; for $\phi\in\ssC^{d,2}(\fg_-;\fg)=\Hom\bigl(\Wedge^2\fg_-;\fg\bigr)_d$ we have
\begin{align*}
	(\partial \phi)(X,Y,Z) &= \comm{X}{\phi(Y,Z)} - \phi(\comm{X}{Y},Z) + (-1)^{|X|(|Y|+|Z|)} \qty( \comm{Y}{\phi(Z,X)} - \phi(\comm{Y}{Z},X) )\\
		& + (-1)^{(|X|+|Y|)|Z|} \qty( \comm{Z}{\phi(X,Y)} - \phi(\comm{Z}{X},Y) ).
\end{align*}

\subsubsection{Filtered deformations}
\label{sec:filtered-def}

A \emph{filtered deformation} of a $\ZZ$-graded Lie superalgebra $\fg$ is a filtered Lie superalgebra $\fgtilde$ whose associated graded superalgebra $\Gr \fgtilde$ is isomorphic to $\fg$ as a~$\ZZ$-graded Lie superalgebra.\footnote{Filtrations of Lie superalgebras are defined in Appendix~\ref{sec:filtered-superspace-superalg}; see Appendix~\ref{sec:filtered-alg} for more on filtrations and associated graded structures.} We will justify this terminology by showing how the bracket of $\fgtilde$ can be viewed explicitly as a~deformation of that of $\fg$ on the same underlying vector space. In parts of this work, we will also say that a filtered Lie superalgebra $\fgtilde'$ is a \emph{filtered subdeformation} of $\fg$ if there exists an embedding of $\ZZ$-graded superalgebras $\Gr\fgtilde'\hookrightarrow \fg$; that is, if $\Gr\fgtilde'\cong\fg'$ as $\ZZ$-graded superalgebras for some graded subalgebra $\fg'$ of $\fg$, although we often do not specify the subalgebra $\fg'$.

Let us fix a filtered Lie superalgebra $\fgtilde=\fgtilde^{-h}\supseteq \fgtilde^{-h+1} \supseteq \cdots $ of finite depth $h$ and choose, \textit{non-canonically}, some complementary subspaces $W_k$ such that $\fgtilde^k=W_k\oplus\fgtilde^{k+1}$ and $W_k\subseteq\fgtilde_{\overline k}$ (which we can demand due to the superspace filtration compatibility conditions \eqref{eq:filter-super-1} and~\eqref{eq:filter-super-2}). As we have already seen, this induces a graded vector superspace isomorphism~${\Phi\colon\fgtilde=\bigoplus_{k=-h}^\infty W_k\to \Gr\fgtilde}$ where for $X\in W_k$, $\Phi(X)=\overline{X}\in\Gr_k\fgtilde$, where the overline denotes projection $\fgtilde^k\to\Gr_k\fgtilde=\fgtilde^k/\fgtilde^{k+1}$.

Now let $X\in W_k$ and $Y\in W_l$. Then since $\fgtilde^{k+l} = W_{k+l}\oplus\fgtilde^{k+l+1} = W_{k+l}\oplus W_{k+l+1}\oplus\cdots$ and the bracket is filtered, we have
\[
	\comm{X}{Y} = \pr_{k+l}\comm{X}{Y} + \pr_{k+l+1}\comm{X}{Y} + \cdots,
\]
where $\pr_j\colon\fgtilde^{k+l}\to W_j$ are the projection maps. We thus define a \emph{graded} bracket $\comm{-}{-}_0$ on $\fgtilde=\bigoplus_{k=-h}^\infty W_k$ as well as a sequence of maps $\bigl(\mu_d\colon\Wedge^2\fgtilde\to\fgtilde\bigr)_{d> 0}$, where $\mu_d$ has degree $d$, by declaring that, on homogeneous elements $X\in W_k$ and $Y\in W_l$,
\[
	\comm{X}{Y}_0 := \pr_{k+l}\comm{X}{Y},	\qquad \mu_d(X,Y) := \pr_{k+l+d}\comm{X}{Y}
\]
and extending linearly. It can be easily verified that $\comm{-}{-}_0$ is a Lie superalgebra bracket; alternatively, it follows from our next observation. Then on homogeneous elements $X\in W_k$ and~${Y\in W_l}$ again, we have
\[
	\comm{\Phi(X)}{\Phi(Y)} = \comm{\overline{X}}{\overline{Y}}_{\mathrm{Gr}} = \overline{\comm{X}{Y}} = \overline{\comm{X}{Y}_0} = \Phi(\comm{X}{Y}_0),
\]
where in each equality we have applied a definition, and in the last equality we have used the fact that $\comm{X}{Y}_0\in W_{k+l}$ is homogenous. Thus, since homogeneous elements span $\fgtilde$, $\Phi$ is a Lie superalgebra isomorphism\footnote{This shows in particular that although the graded bracket $\comm{-}{-}'_0$ defined by a different choice of complementary subspaces is not equal to $\comm{-}{-}_0$, the graded Lie superalgebra structures defined by the two brackets are isomorphic.}
$\bigl(\fgtilde,\comm{-}{-}_0\bigr)\cong \bigl(\Gr\fgtilde,\comm{-}{-}_{\mathrm{Gr}}\bigr)$.

Let us now change perspective slightly. What we have essentially shown here is that if $\fg=\bigoplus_{k=-h}^\infty\fg_k$ is a fixed $\ZZ$-graded superalgebra, a filtered deformation $\fgtilde$ of $\fg$ can be thought of as an alternative Lie algebra structure on the same underlying vector space as $\fg$; now denoting by $\comm{-}{-}$ the graded bracket on $\fg$ and by $\comm{-}{-}'$ the bracket of $\fgtilde$, for $X,Y\in\fg$ we have
\begin{equation}\label{eq:defining-sequence}
	\comm{X}{Y}' = \comm{X}{Y} + \sum_{d>0}\mu_d(X,Y)
\end{equation}
for some positive-degree maps $\bigl(\mu_d\colon\Wedge^2\fg\to\fg\bigr)_{d> 0}$; that is, the primed bracket can be thought of as a deformation of the unprimed bracket by these maps. We will call this sequence a~\emph{defining sequence} of the filtered deformation. This sequence is not unique; as we have already seen, describing the filtered deformation in this way depends on the choices of complementary subspaces. We also cannot simply define a deformed bracket by choosing some arbitrary sequence~$(\mu_d)_{d>0}$; the Jacobi identity for $\comm{-}{-}'$ imposes relations on the defining sequence. Thus, we need a particular algebraic tool to understand the possible filtered deformations of $\fg$, and that is Spencer cohomology.

%%%%%%%%%% {Relation to Spencer cohomology}

\begin{Proposition}\label{prop:def-seq-spencer-cocycle}
	Let $\fg=\bigoplus_{k=-h}^\infty\fg_k$ be a graded Lie superalgebra and $\fgtilde$ a filtered deformation with defining sequence $\bigl(\mu_d\colon\Wedge^2\fg\to\fg\bigr)_{d> 0}$. Then
	\begin{enumerate}\itemsep=0pt
		\item[$(1)$] {\rm\cite[\emph{Proposition}~2.1]{Cheng1998}} The first non-zero term $\mu_k$ in $\qty(\mu_d)_{d> 0}$ is an even Chevalley--Eilenberg $2$-cocycle of $\fg$ with values in the adjoint representation
		$\mu_k\in\ssZ^{k,2}(\fg;\fg)_{\overline{0}}$,	
		where the adjoint Chevalley--Eilenberg complex of $\fg$ is $\ZZ$-graded in a similar fashion to the Spencer complex $($see \eqref{eq:spencer-complex-grading}$)$.
		\item[$(2)$] {\rm\cite[\emph{Proposition}~2.2]{Cheng1998}} Restricting $\mu_k$ to $\fg_-$ gives us a Spencer cocycle
		$\mu_k|_{\Wedge^2\fg_-}\in\ssZ^{k,2}(\fg_-;\fg)$.
		\item[$(3)$] {\rm\cite[\emph{Proposition}~2.2]{Cheng1998}} The cohomology class of this cocycle is $\fg_0$-invariant
			$\qty[\mu_k|_{\Wedge^2\fg_-}]\in \ssH^{k,2}(\fg_-;\fg)^{\fg_0}$.
	\end{enumerate}
\end{Proposition}
We will not prove this or any of the following statements here; we simply note that the above essentially follows by expanding the Jacobi identity for the deformed bracket and treating it degree-by-degree. We will do this for some explicit examples in Section~\ref{sec:filtered-def-poincare}. The results below are more involved but ultimately follow by the same kind of reasoning. We extract from \cite{Cheng1998} only the statements we will need in the rest of this work; see loc.\ cit.\ for full proofs and a more general treatment.

\begin{Definition}\label{def:fund-trans-prolong}
	Let $\fg$ be a graded Lie superalgebra of finite depth $h$. Then $\fg$ is said to be
	\begin{itemize}\itemsep=0pt
	\item \emph{fundamental} if $\fg_-$ is generated by $\fg_{-1}$;
	\item \emph{transitive} if for all $X\in\fg_k$ with $k\geq 0$, $\comm{X}{\fg_-}=0\implies X=0$;
	\item \emph{a full prolongation $\big($of $\bigoplus_{k=-h}^0\fg_k\big)$ of degree $k$} if $\ssH^{d,1}(\fg_-;\fg)=0$ for all $d\geq k$.
	\end{itemize}
\end{Definition}

We will not need fundamentality for the following result, but it is useful to record it here since we will often consider this property alongside the other two defined above.

\begin{Proposition}\label{prop:def-seq-redef}
	Let $\fgtilde$, $\fgtilde'$ be two filtered deformations of $\fg$ with defining sequences
\[
\bigl(\mu_d\colon\Wedge^2\fg\to\fg\bigr)_{d> 0}, \qquad v\bigl(\mu'_d\colon\Wedge^2\fg\to\fg\bigr)_{d> 0},
\]
 respectively. Then
	\begin{enumerate}\itemsep=0pt
		\item[$(1)$] {\rm\cite[\emph{Proposition}~2.3]{Cheng1998}} If for some $k\geq 1$, $\mu_k-\mu_k'|_{\Wedge^2\fg_-}\in\ssB^{k,2}(\fg_-;\fg)$, then $\fgtilde'$ has a defining sequence \smash{$\qty(\mu''_d)_{d> 0}$} such that $\mu''_i=\mu'_i$ for $i<k$ and \smash{$\mu''_k|_{\Wedge^2\fg_-}=\mu_k|_{\Wedge^2\fg_-}$}.
			\label{item:def-seq-redef-1}
		\item[$(2)$] {\rm\cite[\emph{Proposition}~2.5]{Cheng1998}} If $\fg$ is transitive, then the defining sequence $\qty(\mu_d)_{d> 0}$ is completely determined by its restriction to $\fg_-\otimes\fg$.
			\label{item:def-seq-redef-2}
		\item[$(3)$] {\rm\cite[\emph{Proposition}~2.6]{Cheng1998}} If $\fg$ is transitive and for some $k\geq 1$, $\fg$ is an almost full prolongation of degree $k$, $\mu_i|_{\fg_-\otimes\fg}=\mu'_i|_{\fg_-\otimes\fg}$ for $i<k$ and $\mu_k|_{\Wedge^2\fg_-}=\mu'_k|_{\Wedge^2\fg_-}$, then $\fgtilde'$ has a defining sequence \smash{$\qty(\mu''_d)_{d> 0}$} such that $\mu''_i=\mu'_i$ for $i\leq k$ and \smash{$\mu''_i|_{\Wedge^2\fg_-}=\mu_i'|_{\Wedge^2\fg_-}$} for all $i$.
			\label{item:def-seq-redef-3}
	\end{enumerate}
\end{Proposition}

This is a rather technical set of results, but it essentially allows us to redefine the defining sequence of a filtered deformation to partially align it with that of another; under certain conditions, one can completely align the two sequences in order to show that two deformations are isomorphic. We will make use of this in Section \ref{sec:filtered-def-poincare}.

\subsubsection{A note on ``polarised'' and ``depolarised'' expressions}
\label{sec:polarisation}

We finish this section by making note of a useful technique we will often employ in calculations involving Lie superalgebras and elsewhere maps on symmetric products of vector spaces appear.

Let $U$, $W$ be some (finite-dimensional) vector spaces and let $\Psi\in\Odot^n U\to W$ be some completely symmetric multilinear map. Then $\Psi$ is completely determined by ``polarised'' values of the form $\Psi(u,u,\dots,u)$ for elements $u\in U$. The most well-known application of this is the polarisation identity which is used to recover an inner product from its induced norm in the~${n=2}$ case, but it holds more generally. This greatly simplifies calculations involving such maps. When we seek to show that a symmetric tensor vanishes or that two such tensors agree, we will generally do this by showing that the desired property holds in its polarised form, where all inputs are taken to be the same. When we need to ``depolarise'' such an expression, we will indicate that this has been done but will not offer further explanation than what has been indicated here.

\section{Killing spinors and Killing (super)algebras}
\label{sec:killing-spinors-superalgebras}

Let $(M,g)$ be a connected pseudo-Riemannian spin manifold of signature $(p,q)$. We note in particular that $M$ is oriented and denote the special orthonormal frame bundle by $F_{\rm SO}\to M$ and the spin structure by $\varpi\colon P\to F_{\rm SO}$. We need not assume time-orientability,\footnote{Thus $(M,g)$ need not be \emph{strongly spin} in the sense of \cite{Cortes2021,Shahbazi2024_2,Shahbazi2024_1}; see Appendix~\ref{sec:spin-structures} for more details.}
but this assumption may be useful in practice (see Lemma~\ref{lemma:bundle-dirac-current-exist}).
We denote the Lie algebra of Killing vector fields\footnote{Note that this is \emph{not} necessarily the Lie algebra of the isometry group of $(M,g)$ which is in general a subalgebra of $\fiso(M,g)$. We will not need to discuss the isometry group itself so this should not cause confusion.}
by $\fiso(M,g)$ and the Levi-Civita connection by $\nabla$.

Let $S$ be a (possibly $N$-extended) spinor module of $\Spin(p,q)$ and define the \emph{spinor bundle}~${\Sbundle:=P\times_{\Spin(p,q)} S}$; whose sections are the \emph{spinor fields}. We denote the space of spinor fields by $\fS=\Gamma(\Sbundle)$.

We note that the adjoint bundle $\ad F_{\rm SO}\cong\ad P$ can be identified with the bundle of $g$-skew-symmetric endomorphisms of $TM$; we denote its space of sections by $\fso(M,g)$ and recall that it naturally embeds in the Clifford bundle $\Cl(M,g)$ (under natural identification of the latter with the exterior bundle $\Wedge^\bullet TM$) as $\Wedge^2 TM$. As a Lie algebra, the space of sections $\fso(M,g)$ acts on spinor fields via the Clifford action.

See Appendix~\ref{sec:lie-der} for more details of these constructions.

\subsection{Admissible connections and Killing (super)algebras}
\label{sec:admissible-ksa}

Let us now examine the additional structures on the spinor bundle, and the conditions upon them, which are required to define a Killing superalgebra.

\subsubsection{Connections on spinors}

Let $D$ be a connection on the spinor bundle $\Sbundle$. We denote by $\fS_D$ the real vector space of parallel spinors with respect to $D$:
$\fS_D := \qty{\epsilon\in\fS\mid D\epsilon = 0}$.
This is finite-dimensional; indeed, a parallel section can be reconstructed by parallel transport from its value at any point, so in particular we must have $\dim\fS_D\leq\rank\Sbundle=\dim S$. We say that $\fS_D$ is \emph{maximal} if it has the maximal dimension $\dim S$.

Since $\nabla$ and $D$ are both connections on $\Sbundle$, their difference is an $\End\Sbundle$-valued 1-form. We thus define $\beta\in\Omega^1(M;\End\Sbundle)$ by $\beta := \nabla - D$. We now recall that Killing vectors have a natural action on spinors via the \emph{spinorial Lie derivative} (see Appendix~\ref{sec:lie-der-spinors}) and consider whether this action preserves $\fS_D$. For all $X\in\fiso(M,g)$, $Y\in\fX(M)$ and $\epsilon\in\fS_D$, using properties of this Lie derivative (equation~\eqref{eq:lie-der-compat} and a Leibniz rule) for the second equality we have
\begin{align*}
	D_Y\eL_X \epsilon
	&= \nabla_Y\eL_X\epsilon -\beta(Y)\eL_X\epsilon\\
&= \eL_X\nabla_Y\epsilon -\nabla_{\eL_X Y}\epsilon -\eL_X\qty(\beta(Y)\epsilon) +\qty(\eL_X\beta)(Y)\epsilon +\beta(\eL_XY)\epsilon\\
	&= \eL_X D_Y\epsilon - D_{\eL_X Y}\epsilon +\qty(\eL_X\beta)(Y)\epsilon,
\end{align*}
so the spinorial Lie derivative along $X$ preserves $\fS_D$ if and only if $\eL_X\beta$ annihilates $\fS_D$. For the sake of simplicity, and because it will capture the case of supergravity which is our ultimate interest, we will choose to work with Killing vectors which preserve $\beta$; we denote by $\fV_D$ the real subspace of $\fiso(M,g)$ consisting of such vectors
\[
\fV_D := \qty{X\in\fiso(M,g)\mid \eL_X \beta = 0}.
\]
This is in fact an ideal of $\fiso(M,g)$ since it is the annihilator of $\beta$ in the representation of $\fiso(M,g)$ on~$\Omega^1(M;\End\Sbundle)$ defined by the Lie derivative. Clearly, the representation of $\fiso(M,g)$ on~$\fS$ (also via the Lie derivative) restricts to a representation of $\fV_D$ on $\fS_D$.

The curvature of the connection $D$ is the section $R^D \in \Omega^2(M;\End\Sbundle)$ defined by
\begin{gather}\label{eq:D-curvature-def}
	R^D(X,Y)\epsilon = D_XD_Y\epsilon - D_YD_X\epsilon - D_{\comm{X}{Y}}\epsilon,
\end{gather}
where $X,Y\in\fX(M)$ and $\epsilon\in\fS$. One can see immediately that $R^D\epsilon=0$ for all $\epsilon\in\fS_D$, and the former equation is very useful as an integrability condition for the existence of $D$-parallel spinors. In applications, we use the following result to compute $R^D$ and perform manipulations involving it.

\begin{Proposition}\label{prop:D-curvature}
	The curvature $R^D$ is given by the following formula, for $X,Y\in\fX(M)$
	\begin{equation}\label{eq:D-curvature-formula}
		R^D(X,Y)\epsilon = R(X,Y)\cdot\epsilon + \comm{\beta(X)}{\beta(Y)}\epsilon - (\nabla_X\beta)(Y)\epsilon - (\nabla_Y\beta)(X)\epsilon,
	\end{equation}
	where $R$ is the Riemann curvature $($understood as a $2$-form with values in $\ad F_{\rm SO})$ and $[\beta(X),\allowbreak\beta(Y)]$ is a commutator of endomorphisms.
\end{Proposition}

\begin{proof}
	For $\epsilon\in\fS$, we expand $D=\nabla-\beta$ in \eqref{eq:D-curvature-def} to find
	\begin{align*}
		 R^D(X,Y)\epsilon	
			={}& \nabla_X(\nabla_Y\epsilon-\beta(Y)\epsilon) - \beta(X)(\nabla_Y\epsilon-\beta(Y)\epsilon)	\\
				& - \nabla_Y(\nabla_X\epsilon-\beta(X)\epsilon) + \beta(Y)(\nabla_X\epsilon-\beta(X)\epsilon)
				- \nabla_{\comm{X}{Y}}\epsilon + \beta({\comm{X}{Y}})\epsilon	\\
			={}& \qty(\nabla_X\nabla_Y\epsilon - \nabla_Y\nabla_X\epsilon - \nabla_{\comm{X}{Y}}\epsilon)
				+ \qty(\beta(X)\beta(Y)-\beta(Y)\beta(X))\epsilon	\\
				& - \qty(\nabla_X(\beta(Y)\epsilon)-\beta(Y)(\nabla_X\epsilon)) + \qty(\nabla_Y(\beta(X)\epsilon)
				- \beta(X)(\nabla_Y\epsilon)) + \beta({\comm{X}{Y}})\epsilon \\
			={}& R(X,Y)\cdot\epsilon + \comm{\beta(X)}{\beta(Y)} - \nabla_X(\beta(Y))\epsilon + \nabla_Y(\beta(X))\epsilon + \beta({\comm{X}{Y}})\epsilon,
	\end{align*}
	where we have used the definition of the Riemann tensor and the Leibniz rule in the final line. A~further application of the Leibniz rule as well as a the torsion-free property of $\nabla$ gives \eqref{eq:D-curvature-formula}.
\end{proof}

\begin{Remark}\label{rem:riemann-spinor-action}
	It is worth further clarifying the notation $R(X,Y)\cdot\epsilon$ above. As alluded to in the statement of the proposition, the Riemann curvature can be considered as a 2-form with values in $\ad F_{\rm SO}$, meaning that for all $X,Y\in \fX(M)$, $R(X,Y)\in\fso(M,g)$, thus it acts on~$\fS$ via Clifford multiplication. In the proof, we have implicitly used the non-trivial fact that~${R(X,Y)\cdot\epsilon=R^\nabla(X,Y)\epsilon}$ where $R^\nabla\in\Omega^2(M;\End \Sbundle)$ is the curvature 2-form of the spin-lift of the Levi-Civita connection,
\[
R^\nabla(X,Y)\epsilon:=\nabla_X\nabla_Y\epsilon-\nabla_Y\nabla_X\epsilon-\nabla_{\comm{X}{Y}}\epsilon.
\]
Thanks to this relation to the Riemann tensor, we will never have to use this curvature explicitly.
\end{Remark}

\subsubsection{Dirac currents on spinor bundles}
\label{sec:dirac-current-bundle}

In order to define a Killing (super)algebra, we will need to be able to ``square'' (or pair) spinors to tangent vectors, just as we required a Dirac current to define the flat model (super)algebras $\fs$ in Section~\ref{sec:flat-model-algs}. For this purpose, we will define a bundle map $\Otimes^2\Sbundle\to TM$ using a Dirac current on $S$ by exploiting the associated bundle structure of $\Sbundle$ and $TM$ and the holonomy principle.

\begin{Lemma}\label{lemma:bundle-dirac-current-exist}
	Let $V=\RR^{p,q}$ and suppose there exists a Dirac current on the spinor module $S$; that is, an $\fso(p,q)$-equivariant map $\kappa\colon \Odot^2S\to V$ {\rm(}resp.\ $\kappa\colon \Wedge^2S\to V)$. Then there exists a~bundle Dirac current $\underline\kappa\colon \Odot^2\Sbundle\to TM$ {\rm(}resp.\ $\underline\kappa\colon \Wedge^2\Sbundle\to TM)$ which is covariantly constant and which satisfies $\eL_X\underline\kappa=0$ for all $X\in\fiso(M,g)$ if any of the following hold:
	\begin{itemize}\itemsep=0pt
	\item $(M,g)$ has definite signature,
	\item $\kappa$ is $\Spin(p,q)$-equivariant,
	\item $(M,g)$ has indefinite signature and is time-orientable.
	\end{itemize}
\end{Lemma}

\begin{proof}
	We prove the case with $\kappa\colon \Odot^2S\to V$; the other case is completely analogous. Let us consider $\kappa$ as an element $\kappa\in\Odot^2S^*\otimes V$. Fix $x\in M$. We will show that $\kappa$ defines a holonomy-invariant element in the fibre over $x$ of the bundle $\Odot^2\Sbundle^*\otimes TM$ which we note is canonically isomorphic to the associated bundle $E=P\times_{\Spin(p,q)}\bigl(\Odot^2S^*\otimes V\bigr)$. The holonomy principle then guarantees that a covariantly constant section $\underline\kappa$ of that bundle can be defined by parallel transport, which is the desired bundle map.
	
	If $(M,g)$ has definite signature, the spin group is connected, so $\kappa$ is invariant under the action of the spin group, thus the first of our special cases is subsumed by the second. If $(M,g)$ is time-orientable, then the spin structure $\varpi\colon P\to F_{\rm SO}$ is reducible to $\varpi_0\colon P_0\to F_{SO_0}$ (where~$P_0$ is to reduction of $P$ to a $\Spin_0(p,q)$-principal bundle) as discussed in the introduction to Section~\ref{sec:killing-spinors-superalgebras}. In order to treat all three cases in parallel, let us set $G=\Spin(p,q)$ in the first two cases, while in the third case $G=\Spin_0(p,q)$ and $P$ denotes the reduced bundle. In any case, $\kappa$ is $G$-equivariant.
	
	Let us now choose an arbitrary basepoint $u\in P_x$ in the fibre over $x$. This allows us to identify a module of $G$ with the fibre of the associated bundle over $x$; in particular, $\kappa\in\Odot^2S^*\otimes V$ defines an element $\underline\kappa_x:=[u,\kappa]\in E_x$. Moreover, it induces a Lie group isomorphism $\Aut_G(P_x)\cong G$ which allows us to view the holonomy group of the Levi-Civita connection on $P$ at $x$ as a~subgroup of $G$. Thus $G$-invariance of $\kappa$ implies holonomy invariance of $\underline\kappa_x$, and hence the existence of a~covariantly constant global section $\underline\kappa$, as required. We note that $G$-invariance of $\kappa$ also ensures that the definition of $\underline\kappa$ does not depend on the choice of basepoint $u\in P_x$. Finally, $\fg=\fso(p,q)$-invariance of $\kappa$ ensures that $\eL_X\underline\kappa=0$ for all Killing vectors $X$.
\end{proof}

To lighten the notation, we will omit the underline from the bundle map $\kappa$ and trust that this will not cause confusion. With the map $\kappa$ fixed, it will be useful for the proofs of our next few results (and to make a connection to the notation used in Section~\ref{sec:filtered-def-poincare}) to define a section $\gamma$ of the bundle $\Hom\bigl(\Otimes^2\Sbundle,\End(TM)\bigr)$ as follows:
\begin{equation}\label{eq:gamma-geom-def}
	\gamma(\epsilon,\zeta)X := - \kappa(\beta(X)\epsilon,\zeta) - \kappa(\epsilon,\beta(X)\zeta)
\end{equation}
for $\epsilon,\zeta\in\fS$ and $X\in\fX(M)$. Note that $\gamma$ has the same symmetry as $\kappa$. We then have the following.

\begin{Lemma}\label{lemma:dirac-current-endo}
	For all $\epsilon,\zeta\in\fS_D$, $\gamma(\epsilon,\zeta) = A_{\kappa(\epsilon,\zeta)}$.
\end{Lemma}

\begin{proof}
	Let $Z\in\fX(M)$. Then since $\nabla_Z\epsilon = \beta(Z)\epsilon$,
	\begin{gather*}
		A_{\kappa(\epsilon,\zeta)}Z
			= -\nabla_Z\kappa(\epsilon,\zeta)
			= - \kappa(\nabla_Z\epsilon,\zeta) - \kappa(\epsilon,\nabla_Z\zeta)= - \kappa(\beta(Z)\epsilon,\zeta) - \kappa(\epsilon,\beta(Z)\zeta)
			= \gamma(\epsilon,\zeta)
	\end{gather*}
	hence the claim.
\end{proof}

\subsubsection{Existence of the Killing (super)algebra}
\label{sec:exist-ksa}

We now give our first main result followed by a major definition inspired by it.

\begin{Theorem}[existence of (super)algebra associated to $D$]\label{thm:killing-algebra-exist}
	Let $(M,g)$ be a connected pseudo-Riemannian spin manifold, $\Sbundle\to M$ a bundle of spinors with the space of sections $\fS=\Gamma(\Sbundle)$ and let $\kappa\colon \Odot^2\Sbundle\to TM$ {\rm(}resp.\ $\kappa\colon \Wedge^2\Sbundle\to TM)$ be a bundle Dirac current as in Lemma~$\ref{lemma:bundle-dirac-current-exist}$. Let $D$ be a connection on $\Sbundle$ and define $\beta=\nabla-D$, where $\nabla$ is the Levi-Civita spin connection. Define the vector spaces
$\fS_D = \qty{\epsilon\in\fS| D\epsilon = 0}$, $
		\fV_D = \qty{X\in\fiso(M,g)| \eL_X \beta = 0}$.
	and the brackets
$\comm{X}{Y} = \eL_X Y$,
$\comm{X}{\epsilon} = \eL_X \epsilon$,
$\comm{\epsilon}{\zeta} = \kappa(\epsilon,\zeta)$,
	for $X,Y\in\fV_D$, $\epsilon,\zeta\in\fS_D$. Then~${\qty(\fK_D =\fV_D\oplus\fS_D,\comm{-}{-})}$ is a Lie superalgebra $($resp.\ algebra$)$ if and only if the following conditions are satisfied for all $\epsilon,\zeta,\eta\in\fS_D$
	\begin{gather}
		\gamma(\epsilon,\zeta)=A_{\kappa(\epsilon,\zeta)}\in\fso(M,g),\label{eq:killing-alg-exist-1}\\
		\eL_{\kappa(\epsilon,\zeta)}\beta = 0,\label{eq:killing-alg-exist-2}\\
		\beta(\kappa(\epsilon,\zeta))\eta + \beta(\kappa(\zeta,\eta))\epsilon + \beta(\kappa(\eta,\epsilon))\zeta
		+ \gamma(\epsilon,\zeta)\cdot\eta + \gamma(\zeta,\eta)\cdot\epsilon +\gamma(\eta,\epsilon)\cdot\zeta
		= 0 \label{eq:killing-alg-exist-3},
	\end{gather}
	where $\gamma$ is the map defined by equation~\eqref{eq:gamma-geom-def}.
\end{Theorem}

\begin{proof}
	We first check closure and then the Jacobi identities. We have $\comm{\fV_D}{\fV_D}\subseteq\fV_D$ and $\comm{\fV_D}{\fS_D}\subseteq \fS_D$ by construction, while $\comm{\fS_D}{\fS_D}\subseteq\fV_D$ if and only if
	\begin{equation}\label{eq:killing-closure-cond}
		\eL_{\kappa(\epsilon,\zeta)}g = 0
		\qquad\text{and}\qquad
		\eL_{\kappa(\epsilon,\zeta)}\beta = 0 \qquad \forall\epsilon,\zeta\in\fS_D.
	\end{equation}
	By equation~\eqref{eq:killing-condition}, $\eL_{\kappa(\epsilon,\zeta)}g = 0$ if and only if $A_{\kappa(\epsilon,\zeta)}\in\fso(M,g)$. By Lemma~\ref{lemma:dirac-current-endo}, this is equivalent to condition~\eqref{eq:killing-alg-exist-1}. The latter condition in \eqref{eq:killing-closure-cond} is just \eqref{eq:killing-alg-exist-2}.
	
	The Jacobi identity for three vectors is clearly satisfied. For the others, we denote by $[-,-,-]$ the Jacobiator and find
	\[
		[X,Y,\epsilon]= 0 \iff \eL_{\comm{X}{Y}} \epsilon = \eL_X \eL_Y \epsilon - \eL_Y \eL_X \epsilon,
	\]
	which simply says that $\fS_D$ is a representation of $\fV_D$, hence is true by construction;
	\[
		[X,\epsilon,\zeta]= 0 \iff \eL_X \kappa(\epsilon,\zeta) = \kappa(\eL_X\epsilon,\zeta) + \kappa(\epsilon,\eL_X\zeta),
	\]
	which is satisfied by the last part of Lemma~\ref{lemma:bundle-dirac-current-exist};
	\[
		[\epsilon,\zeta,\eta] = 0 \iff \eL_{\kappa(\epsilon,\zeta)} \eta + \eL_{\kappa(\eta,\epsilon)} \zeta + \eL_{\kappa(\zeta,\eta)} \epsilon = 0,
	\]
	which, using the equation $\nabla \epsilon=\beta\epsilon$ and Lemma~\ref{lemma:dirac-current-endo}, is condition~\eqref{eq:killing-alg-exist-3}.
\end{proof}

We interpret this result as giving us a set of conditions on the spinor connection $D$ for the existence of an associated (super)algebra. This inspires the following definition. Note that conditions (1) and (2) are stronger than \eqref{eq:killing-alg-exist-1} and \eqref{eq:killing-alg-exist-3} above; we demand that they hold for \emph{all} spinors, not just $D$-parallel spinors. This makes the definition simpler to check, and it will also be useful for some later results.

\begin{Definition}[admissible connections, Killing spinors and Killing (super)algebras]\label{def:killing-spinor}
	Using the notation of the theorem above, a connection $D=\nabla-\beta$ on $\Sbundle$ is \emph{admissible} if the following hold:
	\begin{enumerate}\itemsep=0pt
		\item The map $\gamma$ defined by equation~\eqref{eq:gamma-geom-def} takes values in $\ad F_{\rm SO}$ (that is, $\gamma(\epsilon,\zeta)\in\fso(M,g)$ for all $\epsilon,\zeta\in\fS$); %\label{item:admiss-conn-1}

		\item $
\beta(\kappa(\epsilon,\zeta))\eta + \beta(\kappa(\zeta,\eta))\epsilon + \beta(\kappa(\eta,\epsilon))\zeta
				+ \gamma(\epsilon,\zeta)\cdot\eta + \gamma(\zeta,\eta)\cdot\epsilon +\gamma(\eta,\epsilon)\cdot\zeta
				= 0$ for all $\epsilon,\zeta,\eta\in\fS$; %\label{item:admiss-conn-2}

		\item $\eL_{\kappa(\epsilon,\zeta)}\beta = 0$ for all $\epsilon,\zeta\in\fS_D$. %\label{item:admiss-conn-3}
	\end{enumerate}
	If $D$ is admissible, the differential equation $D\epsilon=0$ (equivalently $\nabla\epsilon =\beta\epsilon$) is called the \emph{Killing spinor equation} and $\fS_D$ the space of \emph{Killing spinors}. Furthermore, $\fV_D$ is the space of \emph{restricted Killing vectors}, and $\fK_D$ is the \emph{Killing $($super$)$algebra}.\footnote{Our use of the term ``Killing spinor'' here is a generalisation of the sense in which it is used in the supergravity literature and in particular those discussed in \cite{Beckett2021,deMedeiros2016,deMedeiros2018,Figueroa-OFarrill2017_1,Figueroa-OFarrill2017}. It does not subsume all geometric Killing spinors or generalisations thereof; it does however subsume the particular examples in \cite{Figueroa-OFarrill2008_except}.}
\end{Definition}

Let us make some comments to connect the above with treatments of Killing spinors and Killing superalgebras in the supergravity literature. In a supergravity theory, the Killing spinor equation $\nabla\epsilon=\beta\epsilon$ arises by demanding that the supersymmetry variation (with parameter $\epsilon$) of the gravitino vanishes, with the 1-form $\beta$ being parametrised by some bosonic fields $\Phi$. To prove the existence of a Killing superalgebra, one essentially checks that conditions (1) and (2) of the above definition hold and that the Killing spinor equation and bosonic equations of motion imply that $\eL_{\kappa_\epsilon}\Phi=0$, hence $\eL_{\kappa_\epsilon}\beta=0$, where $\kappa_\epsilon$ is the square of a Killing spinor $\epsilon$, which is condition~(3) above.

Often in the supergravity literature, the Killing superalgebra is considered to contain not all of the Killing vectors $X\in\fV_D$, but only those for which $\eL_X\Phi=0$ (hence $\eL_X\beta=0$), giving a subalgebra of what we call $\fK_D$ here in general. Some sources (e.g., \cite{Figueroa-OFarrill2009,Hustler2016}) call this the \emph{symmetry superalgebra}, giving an even more restrictive definition for the Killing superalgebra~-- namely that it contains only those Killing vectors in the span of those generated by pairs of spinors~-- which corresponds to our notion of the \emph{Killing ideal}. Our terminology here agrees with that of~\cite{Figueroa-OFarrill2017_1}. The Killing ideal is referred to as the \emph{transvection ideal} in \cite{Santi2022}.

\begin{Definition}[Killing ideal]
	Let $D$ be an admissible connection and $\fK_D$ its Killing (super)algebra. The \emph{Killing ideal} is the ideal $\overline\fK_D$ of $\fK_D$ generated by the Killing spinors; that is, $\overline\fK_D=\overline\fV_D\oplus\fS_D$ where $\overline\fV_D = \comm{\fS_D}{\fS_D}.$
\end{Definition}

\begin{Remark}
	The conditions \eqref{eq:killing-alg-exist-1}--\eqref{eq:killing-alg-exist-3} (and the related conditions in Definition~\ref{def:killing-spinor}) for existence of a Killing (super)algebra actually concern only the Killing ideal; that is, given a (not necessarily admissible) spinor connection $D$, $\fK_D$ is a Lie (super)algebra if and only if $\overline\fK_D$ is a~Lie (super)algebra. Moreover, we will see later that conditions~(1) and (2) of Definition~\ref{def:killing-spinor} are essentially Spencer cocycle conditions.
\end{Remark}

A similar general formulation of Killing superalgebras to the one given here was given in \cite{Hustler2016}, however our treatment here is more general because we allow for arbitrary metric signature and Dirac current $\kappa$ of arbitrary symmetry, and that work also makes the assumption mentioned above that $\beta$ is parametrised by some bosonic fields $\Phi$ and that the Killing vectors in the algebra satisfy $\eL_X\Phi=0$.

\subsection{Algebraic structure of Killing (super)algebras}
\label{sec:alg-structure-killing}

We now turn to describing the algebraic structure of the (super)algebras $\fK_D$ associated to the admissible connections $D$ we have just defined. Our treatment here is a direct generalisation of that of \cite{Figueroa-OFarrill2017_1}, which considered the case of Killing superalgebras of highly supersymmetric backgrounds of 11-dimensional supergravity. The methods used here were introduced (for Killing vectors) by Kostant \cite{Kostant1955} and independently by Geroch \cite{Geroch1969}, the latter of whom introduced the term \emph{Killing transport}.

\subsubsection{Killing transport}

Let $D$ be an admissible connection on the spinor bundle $\Sbundle$ and recall that the space of Killing spinors ($D$-parallel sections of $\Sbundle$) is denoted $\fS_D$. Let us define the vector bundle
$\eE := TM\oplus\Sbundle\oplus\ad F_{\rm SO}$,
where we recall that the space of sections of $\ad F_{\rm SO}$, the adjoint bundle to the special orthonormal frame bundle, can be identified with the Lie algebra $\fso(M,g)$ of skew-symmetric endomorphisms of the tangent bundle. We define a connection $\eD$ on $\eE$ as follows:
\[
	\eD_Y(X,\epsilon,A) := (\nabla_Y X + AY,D_Y\epsilon,\nabla_Y A + R(Y,X))
\]
for $X,Y\in\fX(M)$, $\epsilon\in\fS$, $A\in\fso(M,g)$.

\begin{Proposition}\label{prop:killing-transport}
	If $D$ is an admissible connection, the parallel sections of $\eD$ are triples $(X,\epsilon,A_X)$, where $X$ is a Killing vector, $\epsilon$ a Killing spinor and $A_X=-\nabla X$. Thus the $\RR$-vector space of $\eD$-parallel sections is isomorphic to $\fiso(M,g)\oplus\fS_D$.
\end{Proposition}
\begin{proof}
If the section $(X,\epsilon,A)$ of $\eE$ is parallel with respect to $\eD$, clearly $\epsilon$ is a Killing spinor and we have $AY = -\nabla_Y X$ for all $Y\in\fX(M)$, so $A = -\nabla X = A_X$. Since $A\in\fso(M,g)$, $X$ must be a Killing vector. Conversely, if $X$ is a Killing vector and $\epsilon$ a Killing spinor, we have $\nabla_Y A_X = -R(Y,X)$ (by equation~\eqref{eq:nabla-A_X}), and so $(X,\epsilon,A_X)$ is parallel with respect to $\eD$.
\end{proof}

Thanks to this result and since $M$ is connected, if one knows the values $(X,\epsilon,A_X)_p$ at any point~${p\in M}$ for any pair $(X,\epsilon)$ with $X\in\fiso(M,g)$ and $\epsilon\in\fS_D$ -- which we call a \emph{Killing pair} -- one can recover the section $(X,\epsilon,A_X)$ by parallel transport with respect to $\eD$. For this reason, $\eD$ is referred to as the \emph{Killing transport connection}, and $(X,\epsilon,A_X)_p$ the \emph{Killing transport data} of~$(X,\epsilon)$ at~$p$. We define an injective $\RR$-linear map $\Phi$ sending a Killing pair in~$\fK_D$ (so that~${X\in\fV_D}$) to its Killing transport data at $p$ as follows:
$\Phi_p\colon \fK_D\to\eE_p$,
	$\Phi_p(X,\epsilon):=(X_p,\epsilon_p,(A_X)_p)$.


\subsubsection{Localising the Killing algebra at a point}

If we restrict our attention to Killing pairs $(X,\epsilon)\in\fK_D$, it is possible to recover the Killing transport data of $\comm{(X,\epsilon)}{(Y,\zeta)}$ at a point from that of $(X,\epsilon)$ and $(Y,\zeta)$. We have
\[
	\comm{(X,\epsilon)}{(Y,\zeta)} = \qty(\comm{X}{Y} + \kappa(\epsilon,\zeta),\eL_X\zeta - \eL_Y\epsilon),
\]
and the corresponding parallel section of $\eE$ is
\begin{equation}\label{eq:killing-comm-triple}
	(\comm{X}{Y} + \kappa(\epsilon,\zeta),\eL_X\zeta - \eL_Y\epsilon,A_{\comm{X}{Y}} + A_{\kappa(\epsilon,\zeta)}).
\end{equation}
We now expand some of the terms in this expression. Equations~\eqref{eq:comm-formula} and \eqref{eq:A-comm} from Appendix~\ref{sec:lei-der-vector-tensor} give us
\[
	\comm{X}{Y} = \nabla_X Y - \nabla_Y X = A_X Y - A_Y X,\qquad
	A_{\comm{X}{Y}} = \comm{A_X}{A_Y} - R(X,Y).
\]
The Killing spinor equation in the form $\nabla \epsilon = \beta\epsilon$ gives
\[
	\eL_Y\epsilon = \nabla_Y\epsilon + A_Y\cdot\epsilon = \beta(Y)\epsilon + A_Y\cdot\epsilon,
\]
and by Lemma~\ref{lemma:dirac-current-endo},
$
	A_{\kappa(\epsilon,\zeta)} = - \kappa(\beta\epsilon,\zeta) - \kappa(\epsilon,\beta\zeta)$.
Thus the triple \eqref{eq:killing-comm-triple} is given by
\begin{gather*}
	\comm{X}{Y} + \kappa(\epsilon,\zeta) = A_X Y - A_Y X + \kappa(\epsilon,\zeta),\\
	\eL_X\zeta - \eL_Y\epsilon = \beta(X)\zeta - \beta(Y)\epsilon + A_X\cdot\zeta - A_Y\cdot\epsilon,\\
	A_{\comm{X}{Y}} + A_{\kappa(\epsilon,\zeta)} = \comm{A_X}{A_Y} - R(X,Y) - \kappa(\beta\epsilon,\zeta) - \kappa(\epsilon,\beta\zeta).
\end{gather*}
We now evaluate the expressions above at $p$ to obtain the Killing transport data of $\comm{(X,\epsilon)}{(Y,\zeta)}$. Let us define the inner product space $(V,\eta):=(T_pM,g_p)$ and identify $S:=\Sbundle_p$ as a ($N$-extended) spinor module of $\Spin(V)\subseteq\Cl(V)$ (where we omit $\eta$ as usual); we then have $\eE_p = V\oplus S\oplus\fso(V)$ as a $\Spin(V)$-module.\footnote{We note that there is some abuse of notation here; we previously denoted the spinor module of $\Spin(p,q)$ used to define $\Sbundle$ by~$S$; identification of these two $S$'s is not canonical and requires e.g., a choice of basepoint in the fibre of the spin structure over~$p$.}
For the sake of readability, we consider the data for the brackets~$\comm{X}{Y}$, $\comm{X}{\epsilon},$ and $\comm{\epsilon}{\zeta}$ separately; we find
\begin{align}
	\Phi_p(\comm{X}{Y})
		&= ((A_X)_p Y_p - (A_Y)_p X_p,0,\comm{(A_X)_p}{(A_Y)_p} - R_p(X_p,Y_p)),
		\label{eq:phi-p-XY}\\
	\Phi_p(\comm{X}{\epsilon})
		&= (0,\beta_p(X_p)\epsilon_p + (A_X)_p\cdot\epsilon_p,0),
		\label{eq:phi-p-Xepsilon}\\
	\Phi_p(\comm{\epsilon}{\zeta})
		&= (\kappa_p(\epsilon_p,\zeta_p),0,- \kappa_p(\beta_p\epsilon_p,\zeta_p) - \kappa_p(\epsilon_p,\beta_p\zeta_p)),
		\label{eq:phi-p-epsilonzeta}
\end{align}
where $\kappa_p\in\Otimes^2S^*\otimes V$, $\beta_p\in V^*\otimes\End S$, and $R_p\in\Wedge^2V^*\otimes \fso(V)$ are the values of $\kappa$, $\beta,$ and $R$ at $p\in M$. Thus we can indeed express the Killing transport data of any $\comm{(X,\epsilon)}{(Y,\zeta)}$ at $p\in M$ in terms of the Killing transport data of $(X,\epsilon)$ and $(Y,\zeta)$ at $p$ if we know the ``background'' data~$(\kappa_p,\beta_p,R_p)$.

\subsubsection{The structure theorem}

We will now prove our second main result. This is a generalisation of the first part of \cite[Theorem~12]{Figueroa-OFarrill2017_1}, which is particular to 11-dimensional supergravity. Recall from Definition~\ref{def:flat-model-alg} that for a fixed inner product space $(V,\eta)$, a flat model (super)algebra $\fs$ is defined by a spinor module $S$ and a Dirac current $\kappa$ on $S$. See Section~\ref{sec:filtered-def} for background on filtered deformations and Appendices~\ref{sec:filtered-alg} and~\ref{sec:filtered-superspace-superalg} for more on filtrations.

\begin{Theorem}[structure of Killing (super)algebras]\label{thm:killing-algebra-filtered}
	Let $D$ be an admissible connection $($see Definition~{\rm\ref{def:killing-spinor})} on a spinor bundle $\Sbundle$ with Dirac current $\kappa$ over a connected pseudo-Riemannian spin manifold $(M,g)$. Then the Killing $($super$)$algebra $\fK_D$ is a filtered deformation of a graded subalgebra $\fa$ of the flat model $($super$)$algebra $\fs$ with odd part $S$ and Dirac current $\kappa$.
\end{Theorem}

\begin{Remark}
Before proving this result, we note that there is a more na\"{\i}ve approach to it which unfortunately fails. We have already observed that $\eE_p = V\oplus S\oplus\fso(V)$, so one might be tempted to identify $\eE_p$ with the flat model algebra $\fs$ and proceed as follows: we have a~map~${\Phi_p\colon \fK_D\to \eE_p=\fs}$, and we can therefore compute brackets on $\Im\Phi_p$,
\begin{gather*}
	\comm{\Phi_p(X)}{\Phi_p(Y)}
		= \qty((A_X)_p Y_p - (A_Y)_p X_p,0,\comm{(A_X)_p}{(A_Y)_p}),\\
	\comm{\Phi_p(X)}{\Phi_p(\epsilon)}
		= \qty(0,(A_X)_p\cdot\epsilon_p,0),\qquad
	\comm{\Phi_p(\epsilon)}{\Phi_p(\zeta)}
		= \qty(\kappa_p(\epsilon_p,\zeta_p),0,0).
\end{gather*}
Comparing to equations~\eqref{eq:phi-p-XY}--\eqref{eq:phi-p-epsilonzeta}, we see that $\Phi_p$ is not a Lie algebra homomorphism in general, but perhaps it is the natural map carrying a filtered Lie algebra to its associated graded inside $\fs$.

The problem is that $\Im\Phi_p$ need not be closed under the bracket $\comm{-}{-}$ -- the expressions above need not be the transport data of any Killing pairs -- and it may not even be a graded subspace of $\eE_p$, since for example Killing vectors embed diagonally; $X\mapsto(X_p,0,(A_X)_p)$. A~more careful method is therefore required.
\end{Remark}

\begin{proof}[Proof of Theorem~\ref{thm:killing-algebra-filtered}]	
	We fix an arbitrary point $p\in M$ and set $V=T_pM$, $S=\Sbundle_p$ as before. Define the evaluation maps
\smash{$\evaluate^V_p\colon \fV_D \to V$}, $X\mapsto X_p$,
\smash{$\evaluate^S_p\colon \fS_D \to S$}, $\epsilon\mapsto \epsilon_p$.
	Let $V'$ be the subspace of $V$ consisting of values of restricted Killing vectors
$V' = \Im\evaluate^V_p = \qty{X_p\in V \mid X\in\fV_D }$,
	and let $\fh$ be the subspace of $\fso(V)$ consisting of values at $p$ of endomorphisms $A_X$ for Killing vectors which vanish at $X$
	\[
		\fh = \qty{(A_X)_p\in\fso(V) \mid X\in\fV_D\colon X_p=0}.
	\]
	Since a Killing vector $X$ is determined by its Killing transport data $(X_p,(A_X)_p)$, there is an isomorphism of vector spaces $\fh\simeq \ker\evaluate^V_p$. For $A\in\fh$, we denote by $X_A$ the corresponding Killing vector in $\ker\evaluate^V_p$. We also define
$S' = \Im\evaluate^S_p = \qty{\epsilon_p\in S \mid \epsilon\in\fS_D}
$
	and note that there exists a vector space isomorphism $S'\simeq \fS_D$ since a Killing spinor is defined by its value at $p$. For ease of notation, when there is no ambiguity we now identify these spaces; the symbol $\epsilon$ will denote both a Killing spinor and its value at $p$.
	
	Now let $\fa = V'\oplus S'\oplus\fh \subseteq \fs$. We claim that this graded subspace is in fact a graded subalgebra of $\fs$. If $A,B\in\fh$, by equation~\eqref{eq:phi-p-XY} we have
	\begin{gather*}
		\comm{X_A}{X_B}_p = 0,\qquad
		(A_{\comm{X_A}{X_B}})_p = \comm{A_{X_A}}{A_{X_B}}_p = \comm{A}{B},
	\end{gather*}
	so $\comm{X_A}{X_B}\in\ker\evaluate^V_p$, and the corresponding element in $\fh$ is $\comm{A}{B}$, showing that $\fh$ is closed under the bracket. Note that this also shows that $\fh\cong\ker\evaluate^V_p$ as Lie algebras. Now let $A\in\fh$ and $v\in V'$. Then there is some $Y\in\fV_D$ with $Y_p=v$. Again by equation~\eqref{eq:phi-p-XY}, we have
	\[
		\comm{X_A}{Y}_p = \qty(A_{X_A})_p Y_p = A v = \comm{A}{v},
	\]
	so $\comm{A}{v}$ is the value at $p$ of $\comm{X_A}{Y}$. Thus $\comm{\fh}{V'}\subseteq V'$. For $A\in\fh$ and $\epsilon\in S'\simeq\fS_D$, by equation~\eqref{eq:phi-p-Xepsilon} we have
	\[
		\comm{X_A}{\epsilon}_p = (A_{X_A})_p\cdot\epsilon_p = A\cdot\epsilon = \comm{A}{\epsilon},
	\]
	so $\comm{A}{\epsilon}\in S'$. Thus $\comm{\fh}{S'}\subseteq S'$. Finally, for $\epsilon,\zeta\in S$, we have
$\comm{\epsilon}{\zeta}_p
			= \kappa_p(\epsilon_p,\zeta_p)
			= \comm{\epsilon_p}{\zeta_p}$,
	so $\comm{S'}{S'}\subseteq V'$. Hence $\fa$ is indeed a graded Lie subalgebra of $\fs$.
	
	Using the notation and terminology of Appendices~\ref{sec:filtered-alg} and~\ref{sec:filtered-superspace-superalg}, we can equip $\fK_D$ with a~filtration $\fK_D^\bullet$ of depth 2 ($\fK_D^i = \fK_D$ for $i\leq-2$) as follows: we set $\fK_D^i = 0$ for $i>0$ and
	\[
		\fK_D^{-2}=\fK_D = \fV_D\oplus\fS_D \quad \supset \quad
		\fK_D^{-1} = \ker\evaluate^V_p\oplus \fS_D \quad \supset \quad
		\fK_D^0 = \ker\evaluate^V_p\cong \fh.
	\]
	The calculations above show that, as $\fh\cong\ker\evaluate^V_p$-modules,
	\begin{gather*}
		\Gr_{-2}\fK_D^\bullet = \faktor{\fK_D}{\ker\evaluate^V_p\oplus\fS_D} \cong V' \qquad \text{\big(using $\evaluate^V_p$\big)},	\\
		\Gr_{-1}\fK_D^\bullet = \faktor{\ker\evaluate^V_p\oplus\fS_D}{\ker\evaluate^V_p} \cong S'	\qquad \text{\big(using $\evaluate^{S}_p$\big)},	\\
		\Gr_0 \fK_D^\bullet = \ker\evaluate^V_p \cong \fh \qquad\text{(using $X\mapsto (A_X)_p$)},
	\end{gather*}
	thus $\Gr\fK_D^\bullet\cong\fa$ as $\fh$-modules. We will show that this isomorphism preserves Lie brackets, hence~$\fK_D$ is a filtered deformation of $\fa$. We denote the associated graded bracket on $\Gr\fK_D^\bullet$ by~$\comm{-}{-}_{\mathrm{Gr}}$ and use an overline to denote projections $\fK_D^i\to\Gr\fK_D^i$. For $A,B\in\fh$, $X,Y\in\fV_D$ and $\epsilon,\zeta\in\fS_D$, we have the following for our map $\Gr\fK_D^\bullet\to\fa$:
	\begin{gather*}
		\comm{X_A}{\overline{Y}}_{\mathrm{Gr}} = \overline{\comm{X_A}{Y}} = \overline{A_{X_A}Y}
			 \longmapsto 	 A Y_p = \comm{A}{Y_p}	\qquad (\text{in degree $-2$}),	\\
		\comm{\epsilon}{\zeta}_{\mathrm{Gr}} = \overline{\comm{\epsilon}{\zeta}} = \overline{\kappa(\epsilon,\zeta)}
			 \longmapsto 	 \kappa_p(\epsilon,\zeta) = \comm{\epsilon}{\zeta}	\qquad (\text{in degree $-2$}),	\\
		\comm{X_A}{\epsilon}_{\mathrm{Gr}} = \overline{\comm{X_A}{\epsilon}} = \overline{A_{X_A}\cdot\epsilon}
			 \longmapsto 	 A\cdot \epsilon = \comm{A}{\epsilon}	\qquad (\text{in degree $-1$}),	\\
		\comm{X_A}{X_B}_{\mathrm{Gr}} = \overline{\comm{X_A}{X_B}} = X_{\comm{A}{B}}
		 \longmapsto 	 \comm{A}{B}	\qquad (\text{in degree $0$}).
	\end{gather*}
	The only brackets we have not considered above are $\big[\overline{X},\epsilon\big]_{\mathrm{Gr}}$ and $\big[\overline{X},\overline{Y}\big]_{\mathrm{Gr}}$, but these lie in $\Gr_{-3}\fK_D^\bullet$ and $\Gr_{-4}\fK_D^\bullet$ respectively, both of which are both zero by definition.
	\end{proof}
	
We recall from the discussion in Section~\ref{sec:filtered-def} that filtered deformations can be explicitly described as a deformation of the bracket of the graded model algebra by a sequence of maps of positive degrees. It will be useful to give our Killing superalgebras such a presentation.

First note that the isomorphism $\fh\cong\ker\evaluate^V_p$ and the evaluation map $\evaluate^V_p$ itself give us the short exact sequence of vector spaces
\[
\begin{tikzcd}
	0 \ar[r] & \fh \ar[r] & \fV_D \ar[r] & V' \ar[r] & 0.
\end{tikzcd}
\]
Any short exact sequence of vector spaces splits, so there exists a (linear) splitting map $\Lambda\colon V'\to \fV_D$ with $\evaluate^V_p(\Lambda(v))=\Lambda(v)_p=v$. We can also parametrise the splitting in terms of a~map $\lambda\colon V'\to \fso(V)$ by setting $\lambda(v)=(A_{\Lambda(v)})_p$; the vector field $\Lambda(v)$ can be recovered from its Killing transport data $(v,\lambda(v))$. The splitting induces a linear isomorphism $V'\oplus\fh\to\fV_D$ given by $(v,A)\mapsto \Lambda(v)+X_A$.
This extends to a $\ZZ_2$-graded linear isomorphism $\Psi\colon \fa \to \fK_D$ given by~${\Psi(v,\epsilon,A) = (\Lambda(v)+X_A,\epsilon)}$,
where $v\in V'$, $A\in\fh$, and $\epsilon\in S'\simeq\fS_D$.
We use $\Psi$ to define a~new bracket $\comm{-}{-}'$ on $\fa$ by pulling back the bracket on $\fK_D$
\[
	\comm{(v,\epsilon,A)}{(w,\zeta,B)}'
		:= \Psi^{-1}(\comm{\Psi(v,\epsilon,A)}{\Psi(w,\zeta,B)})
\]
with respect to which $\Psi$ is a Lie (super)algebra isomorphism $(\fa,\comm{-}{-}')\cong\fK_D$ by construction. Explicitly, the bracket is given by
\begin{gather}
	\comm{A}{B}' = \comm{A}{B},\qquad
	 \comm{A}{v}' = Av + \comm{A}{\lambda(v)}-\lambda(Av),\nonumber	\\
	\comm{v}{w}' = \lambda(v)w-\lambda(w)v + \theta_\lambda(v,w),\qquad
	\comm{A}{\epsilon}' = A\cdot\epsilon,\nonumber	\\
	 \comm{v}{\epsilon}' = \beta_p(v)\epsilon + \lambda(v)\cdot \epsilon,	\qquad
	\comm{\epsilon}{\zeta}'
		 = \kappa_p(\epsilon,\zeta) + \gamma_p(\epsilon,\zeta)-\lambda\qty(\kappa_p(\epsilon,\zeta)),\label{eq:KSA-deformed-brackets}
\end{gather}
where $\gamma_p\colon\Otimes^2S\to\fso(V)$ is the evaluation at $p$ of the map $\gamma$ defined by equation~\eqref{eq:gamma-geom-def} and $\theta_\lambda:\Wedge^2V'\to \fh$ is the map
\[
	\theta_\lambda(v,w)
		= \comm{\lambda(v)}{\lambda(w)} - R_p(v,w) - \lambda(\lambda(v)w-\lambda(w)v).
\]
Using the notation of Lie algebra cohomology, we have a map $\partial\lambda\colon\Wedge^2\fs\to\fs$ of degree $+2$ defined as follows
\begin{gather*}
	\partial\lambda(A,v) = \comm{A}{\lambda(v)}-\lambda(Av),\qquad
	 \partial\lambda(v,w) = \lambda(v)w-\lambda(w)v,\qquad
	\partial\lambda(v,\epsilon)	= \lambda(v)\cdot\epsilon,\\
	 \partial\lambda(\epsilon,\zeta)	= -\lambda(\kappa_p(\epsilon,\zeta)),
\end{gather*}
whence we have
\[
	\comm{-}{-}' = \comm{-}{-} + \beta_p + \gamma_p + \partial\lambda + \theta_\lambda,
\]
where $\beta_p$, $\gamma_p$ and $\theta_\lambda$ have been trivially extended to maps $\Wedge^2\fs\to\fs$. The deforming maps are $\mu = \beta_p + \gamma_p+\partial\lambda$, of degree $+2$, and $\theta_\lambda$, of degree $+4$. That is, in the language of Section~\ref{sec:filtered-def} (see equation \eqref{eq:defining-sequence}), the deformation has defining sequence $(\mu,\theta_\lambda,0,\dots)$.

The splitting $\Lambda\colon V'\to \fV_D$ and the corresponding map $\lambda \colon V'\to\fh$ can be viewed as a ``correction'' to the failed na\"{\i}ve approach discussed in the remark before the proof of Theorem~\ref{thm:killing-algebra-filtered}. Note that these maps are not canonical; if $\Lambda'\colon V'\to\fV_D$ is another splitting and $\lambda'\colon V'\to\fso(V)$ is the corresponding map, we have
\[
	\Phi_p(\Lambda'(v)-\Lambda(v)) = (v-v,\lambda'(v)-\lambda(v)) = (0,\lambda'(v)-\lambda(v)),
\]
so $(\Lambda'-\Lambda)(v)\in\ker\evaluate^V_p$, and $(\lambda'-\lambda)(v)$ is the corresponding element of $\fh$. Thus $\lambda$, and therefore the splitting, is unique up to a choice of map $V'\to\fh$. We will see later that such a~map is a Spencer $(2,1)$-cocycle for $\fa$, so it is consistent that changing the defining sequence by the coboundary of such a map does not change the isomorphism class of the deformation, only its presentation. Also note that this choice of splitting is exactly a choice of complementary submodules as discussed near the beginning of Section~\ref{sec:filtered-def}.

\subsection{Constrained Killing spinors}
\label{sec:alg-kse}

So far in this section, we have only considered the differential Killing spinor equation ${D\epsilon=0}$ which in supergravity theories arises from demanding that the variation of the gravitino ${\partial_\epsilon\Psi=\!D\epsilon}$ vanishes. While in some pure minimal supergravity theories (in particular, those in dimensions~4,~5,~6 and~$11$, which have been studied using Spencer cohomology \cite{Beckett2021,deMedeiros2016,deMedeiros2018,Figueroa-OFarrill2016,Figueroa-OFarrill2017}) this is the only local condition imposed by supersymmetry, in most supergravity theories (including any extended, gauged or matter-coupled theories and even some pure minimal theories), there are additional fermions present. A background of such a theory is called supersymmetric if the variations of \emph{all} fermions vanish for some supersymmetry parameter. The variations of the additional fermions are always algebraic, rather than differential, in the spinor. The vanishing of these variations are referred to as \emph{algebraic} or \emph{linear constraints} and we will call the spinors satisfying them as well as $D\epsilon=0$ \emph{constrained Killing spinors}.\footnote{Our terminology follows \cite{Cortes2021,Lazaroiu2016} (omitting the word ``generalised''). The more recent work \cite{Shahbazi2024_1} uses the term ``constrained differential spinor''.}
Similar considerations for Killing superalgebras are found in \cite{Hustler2016}.

We will make some comments here about how these might be brought into our present framework; we leave a full treatment for future work.

\subsubsection{Linear constraints}

Schematically, the linear constraints are of the form $\eP\epsilon=0$ where $\eP$ is some fibrewise linear operator on the spinor bundle $\Sbundle$. More concretely, all fermions may take values in the same representation as the gravitini, in which case $\eP\in\End\Sbundle$. In a more general situation, the additional fermions take values in some vector bundle $E\to M$ (for example, an associated bundle for some gauge symmetry). Then the constraint equation is $\eP\epsilon=0$, where
\[
\eP\in\Hom(\Sbundle,\Sbundle\otimes E)\cong \End\Sbundle\otimes E.
\]
This bundle map is parametrised by the bosonic fields of the background, just as~${\beta=D-\nabla}$~is.\footnote{Let us note that the fermion variations $\eP\epsilon$ and the expression $\beta\epsilon$ are fibrewise linear in $\epsilon$ but not necessarily in the bosonic background fields. For example, gauge field strengths (hence derivatives of gauge connections) and derivatives of scalars typically appear as coefficients in both, and may also appear in nonlinear combinations.}

Let $D=\nabla-\beta$ be a spinor connection which is not necessarily admissible in the sense of Definition~\ref{def:killing-spinor} and let us denote the space of $D$-parallel sections of $\Sbundle$ which are annihilated by~$\eP$~by
\[
	\fS_{D,\eP} := \qty{\epsilon\in\fS\mid D\epsilon = 0 \text{ and } \eP\epsilon = 0}.
\]
We previously discussed the integrability condition $R^D\epsilon=0$ for the existence of $D$-parallel spinors (see equation~\eqref{eq:D-curvature-def} and the subsequent discussion). In the case of constrained Killing spinors $\epsilon\in\fS_{D,\eP}$, this can be supplemented by further integrability conditions, the first-order such conditions being schematically (see, e.g., \cite{Gran2019})
$\comm{\eP}{\eP}\epsilon = 0$, \smash{$\widehat{D}\eP \epsilon := \comm{D}{\eP}\epsilon = 0$},
where~${\comm{-}{-}}$ in both expressions is a commutator. These conditions require some interpretation which depends on the structure on the bundle $E$ in which the additional fermions take their values. If $\eP\in\End\Sbundle$, so $E\to M$ is the trivial line bundle, then the meaning of the expressions is clear. The same holds if $E$ is simply the sum of such bundles. In the more general situation, we may interpret $\eP$ as acting on $\Sbundle\otimes E$ as $\eP\otimes\Id_E$ and extend $D$ to a connection on $\Sbundle\otimes E$ which acts trivially on the second factor. In practice though, there may be some more natural interpretation in which these operations are ``twisted'' by some structure on $E$.

\subsubsection{Constrained Killing superalgebra}

As we did for $\fS_D$ in Section~\ref{sec:admissible-ksa}, we now seek to find a subalgebra of the algebra of Killing vectors which preserves this space under the action of the Lie derivative; a natural choice is
\[
	\fV_{D,\eP} := \big\{X\in\fiso(M,g) \mid \eL_X \beta = 0 \text{ and } \eLhat_X \eP = 0\big\},
\]
where again we must interpret the expression $\eLhat_X \eP$. In particular, we assume that there is some structure on $E$ such which gives rise to some generalisation $\eLhat$ of the Lie derivative which acts on $\Hom\Sbundle\otimes E$ -- if $E$ happens to be a trivial bundle, the spinorial Lie derivative naturally generalises, while if it is an associated bundle to the principal bundle for some gauge symmetry, a covariant derivative constructed from a gauge connection will suffice. We note that in the latter case, $\fV_{D,\eP}$ may not close under the Lie bracket, so we may need to refine its definition.

If $\fV_{D,\eP}$ (or some more carefully chosen subspace) is indeed a subalgebra of $\fiso(M,g)$, its action via the spinorial Lie derivative preserves $\fS_{D,\eP}$.
We now seek to give $\fK_{D,\eP} = \fV_{D,\eP}\oplus\fS_{D,\eP}$ the structure of a superalgebra using the Lie derivative and Dirac current as in Theorem~\ref{thm:killing-algebra-exist}. But as in the proof of that theorem, all that needs to be checked is that the $\comm{\fS_{D,\eP}}{\fS_{D,\eP}}$ bracket closes on $\fV_{D,\eP}$, and that the $\qty[\fS_{D,\eP},\fS_{D,\eP},\fS_{D,\eP}]$ Jacobi identity is satisfied. But we quickly see that this is the case if equations \eqref{eq:killing-alg-exist-1}--\eqref{eq:killing-alg-exist-3} as well as
\begin{equation}\label{eq:alg-constraint-cons}
	\eLhat_{\kappa(\epsilon,\zeta)}\eP = 0
\end{equation}
are satisfied for all $\epsilon,\zeta\in\fS_{D,\eP}$. We note that in a supergravity theory, $\beta$ and $\eP$ are defined in terms of background fields which we expect to be preserved by the Lie derivatives along Killing vectors generated by spinors in $\fS_{D,\eP}$, hence the condition above is also satisfied.

Assuming that the expressions above can be made precise as suggested, we should then generalise Definition~\ref{def:killing-spinor} to the constrained case, considering not just admissible connections $D$ but admissible \emph{pairs} $(D,\eP)$, incorporating the condition \eqref{eq:alg-constraint-cons} and possibly demanding that the admissibility conditions of Definition~\ref{def:killing-spinor} should only be imposed on $\ker\eP$, rather than the whole spinor bundle.

\subsubsection{Killing transport and filtered deformation}

We now consider whether an analogue of Theorem~\ref{thm:killing-algebra-filtered} holds, that is, whether it is the case that the Killing superalgebra is a filtered deformation of the Poincar\'e superalgebra in the constrained case. Recall that an important part of this proof was the localisation of the Killing vectors and spinors in $\fK_D$ at a point via Killing transport. Since constrained Killing spinors are still $D$-parallel with respect to a connection $D$, the same arguments hold, the only modification being that the Killing transport data of a constrained Killing spinor at $p\in M$ must lie in $\ker\eP_p$. Note that this does not mean that any spinor at $p$ annihilated by $\eP$ can be Killing-transported to define a constrained Killing spinor; as in the unconstrained case, a $D$-parallel global section with this transport data may not exist, and in the constrained case, even if one does exist, it need not be $\eP$-parallel.

One should thus expect an analogue to Theorem~\ref{thm:killing-algebra-filtered} for $\fK_{D,\eP}$ with essentially the same proof, the main difference being that the filtered subdeformation of $\fs$ one thus obtains must have $S'\simeq \evaluate^S_p\fS_{D,\eP}\subseteq\ker\eP_p\subseteq\Sbundle_p$.

We will not consider constrained Killing spinors in the rest of this work, primarily because, in contrast to the differential Killing spinor equation, the linear constraints apparently do not have an interpretation in terms of Lie algebra cohomology. We thus leave the details of the formalism suggested above to be fully treated in future publications.

\section{Filtered subdeformations of the Poincar\'e superalgebra}
\label{sec:filtered-def-poincare}

Theorem~\ref{thm:killing-algebra-filtered} establishes that any Killing (super)algebra is a filtered deformation of a graded subalgebra of a flat model (super)algebra $\fs$ (see Definition~\ref{def:flat-model-alg}). Let us now consider such deformations from a homological point of view. Throughout, let $(V,\eta)$ be an inner product space and let $S$ be a (possibly $N$-extended) spinor module of $\Spin(V)$. For ease of exposition, we now specialise to the case of a symmetric Dirac current $\kappa\colon \Odot^2 S \to V$, whence we work with \emph{super}algebras; the skew-symmetric case is entirely analogous. Eventually, we will further specialise to Lorentzian signature, hence Poincar\'e superalgebras, and make some further assumptions on the Dirac current and the graded subalgebras we work with in order to take advantage of some technical simplifications offered by the homogeneity theorem (see Theorem~\ref{thm:homogeneity}). We comment on how dependent our results are on these assumptions in Remark~\ref{rem:relaxing-homogeneity-assumption}. We will make extensive use of the background material on Spencer cohomology and filtered deformations presented in Section~\ref{sec:spencer-thy}.

\subsection{The Spencer (2,2)-cohomology}
\label{sec:poincare-spencer-22-cohomology}

We saw in Section~\ref{sec:filtered-def} that filtered deformations of graded Lie superalgebras are governed by their Spencer cohomology. We therefore need to understand the Spencer cohomology of the flat model superalgebra and its graded subalgebras before discussing their deformations.

\subsubsection{The complex}

A graded subalgebra $\fa$ of the flat model superalgebra $\fs$ defined by the data $(V,S,\kappa)$ takes the form $\fa=\fa_{-2}\oplus\fa_{-1}\oplus\fa_0$, where $\fa_{-1}=S'$ is a vector subspace of $S$, $\fa_{-2}=V'$ is a subspace of $V$ containing $\Im \kappa|_{\Odot^2S'}$ and $\fa_0=\fh$ is a subalgebra of $\fso(V)$ preserving $S'$ and $V'$. By Proposition~\ref{prop:def-seq-spencer-cocycle}, to study filtered deformations of $\fa$ we must first understand the degree-2 Spencer complex $\bigl(\ssC^{2,\bullet}(\fa_-;\fa),\partial\bigr)$. Recalling the definition of the Spencer complex from Section~\ref{subsec:spencer-cohomology}, the space of $(2,p)$-cochains consists of degree-2 maps $\Wedge^p\fa\to\fa$, whence we deduce that the degree-2 complex can be written as
\begin{align*}
	0 &\longrightarrow \Hom(V',\fh)
\longrightarrow \Hom\bigl(\Wedge^2V',V'\bigr)\oplus\Hom(V'\otimes S',S')\oplus\Hom\bigl(\Odot^2S',\fh\bigr)\\
	&\longrightarrow \Hom\bigl(V'\otimes\Odot^2S',V'\bigr)\oplus\Hom\bigl(\Odot^3S',S'\bigr)\longrightarrow 0.
\end{align*}
We denote the projections to the components of $\ssC^{2,2}(\fa_-;\fa)$ as follows
\begin{gather*}
	\pi_1\colon\ \ssC^{2,2}(\fa_-;\fa) \longrightarrow \Hom\bigl(\Wedge^2V',V'\bigr),\qquad
	\pi_2\colon\ \ssC^{2,2}(\fa_-;\fa) \longrightarrow \Hom(V\otimes S',S'),\\
	\pi_3\colon\ \ssC^{2,2}(\fa_-;\fa) \longrightarrow \Hom\bigl(\Odot^2S',\fh\bigr).
\end{gather*}

We now record a pair of results which will be useful for studying the cohomology of a class of graded subalgebras which we will be particularly interested in later on.

\begin{Lemma}\label{lemma:V-h-maps}
	Let $\fh$ be a subalgebra of $\fso(V)$ and $V'$ an $\fh$-submodule of $V$ such that
	\begin{enumerate}\itemsep=0pt
		\item[$(1)$] the restriction of $\eta$ to $V'$ is non-degenerate,
		\item[$(2)$] the action of $\fh$ on $V'$ is faithful,
	\end{enumerate}
	$($which hold in particular for $V'=V)$ and suppose that $\lambda\colon V'\to\fh$ is a linear map such that
	\[%\label{eq:lambda-symm}
		\lambda(v)w - \lambda(w)v = 0, \qquad \forall v,w\in V'.
	\]
	Then $\lambda=0$.
\end{Lemma}

\begin{proof}
	Using that $\lambda$ takes values in $\fh\subseteq \fso(V)$, for all $u,v,w\in V'$ we have
	\begin{align*}
		\eta(u,\lambda(v)w)
		&= \eta(u,\lambda(w)v)
		= -\eta(\lambda(w)u,v)
		= -\eta(\lambda(u)w,v) \\
		&
		= \eta(w,\lambda(u)v)
		= \eta(w,\lambda(v)u)
		= -\eta(\lambda(v)w,u).
	\end{align*}
	Thus by non-degeneracy we have $\lambda(v)w=0$ for all $v,w\in V'$, so since $\fh$ acts faithfully on $V'$, $\lambda=0$.
\end{proof}

\begin{Corollary}\label{coro:21-22-injective}
	Let $\fa=V\oplus S'\oplus \fh$ be a graded subalgebra of $\fs$ with $\fa_{-2}=V$. Then the composition
$\pi_1\circ\partial\colon \Hom(V,\fh) \to \Hom\bigl(\Wedge^2 V,V\bigr)
$
	is injective, whence ${\ssH^{2,1}(\fa_-;\fa)=\ssZ^{2,1}(\fa_-;\fa)=0}$. If~${\fh=\fso(V)}$, $\pi_1\circ\partial$ is an isomorphism.
\end{Corollary}
\begin{proof}
Let $\lambda\in\ker(\pi_1\circ\partial)$. Then for all $v,w\in V$,
$(\pi_1\circ\partial)(\lambda)(v,w) = \partial\lambda(v,w) = \lambda(v)w - \lambda(w)v = 0$.
The first claim then follows from Lemma~\ref{lemma:V-h-maps}. The second claim follows by a~dimension count.
\end{proof}

\begin{Remark}
	The second claim in the result above can be restated as follows: for any linear map~${\alpha\colon\Wedge^2 V\to V}$, there exists a unique linear map $\lambda\colon V\to\fso(V)$ satisfying $\alpha(v,w)=\partial\lambda(v,w) = \lambda(v)w-\lambda(w)v$. We will use this fact repeatedly.
\end{Remark}

\subsubsection{The cohomology}

The discussion above gives us a characterisation of the $(2,2)$-cohomology group of the whole flat model superalgebra $\fs$ which we will later be able to bootstrap in order to study a particular class of filtered subdeformations.

\begin{Lemma} \label{lemma:H22-full-flat-model}
	Let $\alpha+\beta+\gamma$ be a cocycle in $\ssC^{2,2}(\fs_-;\fs)$ where
	\begin{align}\label{eq:22-cocycle-comps}
		\alpha\in\Hom\bigl(\Wedge^2V,V\bigr),
		\qquad\beta\in\Hom(V\otimes S,S),
		\qquad\gamma\in\Hom\bigl(\Odot^2S,\fso(V)\bigr).
	\end{align}
	Then the cohomology class $[\alpha+\beta+\gamma]\in \ssH^{2,2}(\fs_-;\fs)$ has a unique representative $\beta'+\gamma'$ with no $\Hom\bigl(\Wedge^2 V,V\bigr)$ component. We denote the space of such \emph{normalised cocycles} by $\cH^{2,2}$
	\[
		\cH^{2,2}
		= \qty{\beta+\gamma \in \ssZ^{2,2}(\fs_-;\fs)
			 \mid
				\beta\in\Hom(V\otimes S,S), \gamma\in\Hom\big(\Odot^2S,\fso(V)\big)
			},
	\]
	and then we have
$\ssH^{2,2}(\fs_-;\fs) \cong \cH^{2,2}
$
	as $\fso(V)$-modules.
\end{Lemma}

\begin{proof}
	Using the isomorphism $\pi_1\circ\partial$, there exists a unique map $\lambda\in\Hom(V,\fso(V))$ such that~${\partial\lambda(v,w)=\lambda(v)w - \lambda(w)v = \alpha(v,w)}$.
	Thus, if we define $\beta'$ and $\gamma'$ by
	\begin{gather*}%\label{eq:H22-normalised-def}
		\beta'(v,s) = \beta(v,s) - \partial\lambda(v,s)
			= \beta(v,s) - \lambda(v)\cdot s,\\
		\gamma'(s,s') = \gamma(s,s') - \partial\lambda(s,s')
			= \gamma(s,s) + \lambda(\kappa(s,s')),
	\end{gather*}
	we have
$[\alpha+\beta+\gamma]
		= [\alpha+\beta+\gamma-\partial\lambda]
		= [\beta' + \gamma']$,
	which proves the first claim. It immediately follows that, as $\fso(V)$-modules, we have $\ssZ^{2,2}(\fs_-;\fs)=\cH^{2,2}\oplus\ssB^{2,2}(\fs_-;\fs)$ and $\ssH^{2,2}(\fs_-;\fs) \cong \cH^{2,2}$.
\end{proof}

It is worth pausing to explicitly examine the ``normalised'' Spencer cocycle conditions for $\beta+\gamma\in\cH^{2,2}$. Written in polarised form (explained in Section~\ref{sec:polarisation}), these are
\begin{gather}
	2\kappa(s,\beta(v,s)) + \gamma(s,s)v = 0, \label{eq:cocycle-1}	\\
	\beta(\kappa_s,s) + \gamma(s,s)\cdot s = 0	\label{eq:cocycle-2}
\end{gather}
for $v\in V$ and $s\in S$, where we introduce the notation $\kappa_s:=\kappa(s,s)$. We will refer to these as the \emph{first} and \emph{second} \emph{$($normalised$)$ Spencer cocycle conditions} for $\fs$ respectively. We note that the first equation uniquely defines $\gamma\colon\Odot^2S\to\fso(V)$ in terms of $\beta\colon V\otimes S\to S$ and constrains the endomorphism $v\mapsto \kappa(s,\beta(v,s))$ of $V$ to lie in $\fso(V)$ for all $s\in S$. In explicit cohomology calculations (see \cite{Beckett2021,deMedeiros2016,deMedeiros2018,Figueroa-OFarrill2017}), we first determine the space of maps $\beta$ satisfying this constraint and then substitute the resulting expressions for $\beta$ and $\gamma$ into the second equation and solve it.

We also note here that $\beta+\gamma\in\cH^{2,2}$ (or the corresponding cohomology class) is invariant under the action of a subalgebra $\fh\in\fso(V)$ if and only if $\beta$ is $\fh$-invariant (that is, $\beta$ is $\fh$-equivariant when considered as a map $V\otimes S\to S$); indeed, if $A\in\fh$, $v\in V$ and $s\in S$ then we have~${(A\cdot(\beta+\gamma))(v,s) = (A\cdot\beta)(v,s) = A\cdot (\beta(v,s)) - \beta(Av,s) - \beta(A,v\cdot s)}
$
so clearly~$\beta$ is invariant if $\beta+\gamma$ is, while on the other hand if $\beta$ is invariant then using the first cocycle condition (once in the form above and once in depolarised form) as well as $\fso(V)$-invariance of~$\kappa$, we have\looseness=-1
\begin{align*}
(A\cdot\gamma)(s,s)v &= \comm{A}{\gamma(s,s)}v - 2\gamma(A\cdot s,s)v \\
	&= 2A(\kappa(s,\beta(v,s))) - 2\kappa(s,\beta(Av,s)) - 2\kappa(A\cdot s,\beta(v,s)) - 2\kappa(s,\beta(v,A\cdot s)) \\
	&= 2\kappa(s,A\cdot\beta(v,s)) - 2\kappa(s,\beta(Av,s)) - 2\kappa(s,\beta(v,A\cdot s))= 0,
\end{align*}
thus $\gamma$ is $\fh$-invariant, whence $\beta+\gamma$ is too.

\begin{Remark}
	Let us pause briefly to relate these cocycle conditions back to our theory of Killing superalgebras and to explain the coincidences in notation. First note that there is a canonical homomorphism $\Hom(V\otimes S,S)\cong V^*\otimes\End S$, so we can view the cocycle component $\beta$ as a~spinor endomorphism-valued 1-form on $V$. Under this interpretation, we write $\beta(v)s:=\beta(v,s)$.
	
	The definition of admissible connections, Definition~\ref{def:killing-spinor}, involves three conditions on the 1-form $\beta=D-\nabla\in\Omega^1(M;\End\Sbundle)$ and the section $\gamma\in\Gamma\bigl(\Odot^2\Sbundle^*\otimes TM\bigr)$ which is defined in terms of $\beta$ by equation~\eqref{eq:gamma-geom-def}. Conditions~(1) and~(2) are algebraic in $\beta$, $\gamma$ and condition~(3) is differential. For a symmetric Dirac current, we can polarise the definition of $\gamma$ and the algebraic conditions. Then, taking them pointwise, the definition of $\gamma$ and condition~(1) are equivalent to cocycle condition \eqref{eq:cocycle-1} above, while condition~(2) is precisely \eqref{eq:cocycle-2} (where the fibre of $\Sbundle$ is~$S$ and the same Dirac current is chosen). Thus, if we wish to find admissible connections on a~spinor bundle $\Sbundle$ equipped with a Dirac current $\kappa$ over a spin manifold $(M,g)$, we need only compute~$\cH^{2,2}$ to determine a local ansatz for the allowed 1-forms $\beta$ and then check whether or under what circumstances the remaining condition (namely, $\eL_{\kappa_\epsilon}\beta=0$ for all $D$-parallel spinors~$\epsilon$) holds, which we typically do using integrability conditions.
\end{Remark}

\subsection{General filtered deformations and cohomology}
\label{sec:filtered-def-cohom}

Recall from Section~\ref{sec:filtered-def} that we can describe filtered deformations of a $\ZZ$-graded superalgebra as deformed brackets on the same underlying vector space as the original algebra. In particular, a filtered deformation $\fatilde$ of a graded subalgebra $\fa=V'\oplus S'\oplus \fh$ of $\fs$ has the brackets
\begin{gather}
	\comm{A}{B} = AB-BA,\qquad
	\comm{s}{s} = \kappa_s + \gamma(s,s),\qquad
	\comm{A}{v} = Av + \delta(A,v),
	\nonumber\\
	\comm{v}{s} = \beta(v,s),\qquad\comm{A}{s} = A\cdot s
	, \qquad\comm{v}{w} = \alpha(v,w) + \theta(v,w),\label{eq:def-brackets-general}
\end{gather}
where $A,B\in\fh$, $v,w\in V'$, $s\in S'$, and
$\alpha\colon \Wedge^2V' \to V'$, $
	\beta\colon V'\otimes S' \to S'$, $
	\gamma\colon \Odot^2 S' \to \fh$, $
	\delta\colon \fh\otimes V' \to \fh
$
are the components of the degree-2 deformation map, while
$\theta\colon \Wedge^2 V' \to \fh
$
is the degree-4 deformation map. We denote the full degree-2 deformation by
${\mu=\alpha+\beta+\gamma+\delta\colon \fa\otimes\fa \to \fa}
$
so the defining sequence (as in equation~\eqref{eq:defining-sequence}) for this deformation is $(\mu,\theta,0,\dots)$. From Proposition~\ref{prop:def-seq-spencer-cocycle}, we have the following homological information:
\begin{gather}
	\mu \in \ssZ^{2,2}(\fa;\fa),	\label{eq:mu-CE-cocycle}	\\
	\mu|_{\fa_-\otimes\fa_-} \in \ssZ^{2,2}(\fa_-;\fa),	\label{eq:mu-spencer-cocycle}	\\
	\qty[\mu|_{\fa_-\otimes\fa_-}] \in \ssH^{2,2}(\fa_-;\fa)^\fh, \label{eq:mu-invt-class}
\end{gather}
where in the first expression we have graded the full Chevalley--Eilenberg complex in the same way that we graded the Spencer complex in equation~\eqref{eq:spencer-complex-grading}.

\subsubsection{Unpacking the cocycle conditions}

Let us consider these conditions in more detail. Equation~\eqref{eq:mu-CE-cocycle} is equivalent to the following system of linear equations:
\begin{gather}
	\alpha(\kappa_s,v) + 2\kappa(s,\beta(v,s)) + \gamma(s,s)v = 0,\label{eq:22-spencer-cocycle-vss} \\
	\beta(\kappa_s,s) + \gamma(s,s)\cdot s = 0,\label{eq:22-spencer-cocycle-sss}	\\
	A\alpha(v,w) - \alpha(Av,w) - \alpha(v,Aw) + \delta(A,w)v - \delta(A,v)w = 0,\label{eq:alpha-delta} \\
	A\cdot(\beta(v,s)) - \beta(Av,s) - \beta(v,A\cdot s) - \delta(A,v)\cdot s = 0,\label{eq:beta-delta} \\
	\comm{A}{\gamma(s,s)} - 2\gamma(A\cdot s,s) + \delta(A,\kappa_s) = 0,\label{eq:gamma-delta} \\
	\delta(\comm{A}{B},v) - \comm{A}{\delta(B,v)} + \comm{B}{\delta(A,v)} - \delta(A,Bv) + \delta(B,Av) = 0 \label{eq:delta-coycle-cond}
\end{gather}
for all $A,B\in\fh$, $v,w\in V'$ and $s\in S'$, where again $\kappa_s=\kappa(s,s)$. Equation~\eqref{eq:mu-spencer-cocycle} is equivalent to the first pair of equations in this system -- we will refer to these as the \emph{Spencer cocycle conditions} for $\fa$. Equation~\eqref{eq:mu-invt-class} is equivalent to the following: for each $A\in\fh$, there exists a~cochain $\chi_A\in \ssC^{2,1}(\fa_-;\fa)=\Hom(V',\fh)$ such that
\begin{equation}\label{eq:chiA-def}
	A\cdot(\alpha+\beta+\gamma) = \partial\chi_A.
\end{equation}
This is implied by equation~\eqref{eq:mu-CE-cocycle}: simply taking $\chi_A=\imath_A\delta$, the above equation is equivalent to equations~\eqref{eq:alpha-delta}--\eqref{eq:gamma-delta}. In the case that $V'=V$, we have a converse: it follows from Corollary~\ref{coro:21-22-injective} that \eqref{eq:chiA-def} defines $\chi_A$ uniquely and that the assignment $A\mapsto\chi_A$ is linear. Thus we can define a unique linear map $\delta\colon \fh\otimes V\to\fh$ by $\delta(A,v):=\chi_A(v)$. Then
\begin{gather*}
	\delta(\comm{A}{B},v) - \comm{A}{\delta(B,v)} + \comm{B}{\delta(A,v)} - \delta(A,Bv) + \delta(B,Av)\\
	\qquad= \chi_{\comm{A}{B}}(v) - \comm{A}{\chi_B(v)} + \comm{B}{\chi_A(v)}	- \chi_A(Bv) + \chi_B(Av)	\\
	\qquad= \qty(\chi_{\comm{A}{B}} - A\cdot\chi_B + B\cdot\chi_A)(v).
\end{gather*}
On the other hand, by \eqref{eq:chiA-def} we have
$\partial(\chi_{\comm{A}{B}} - A\cdot\chi_B + B\cdot\chi_A) = 0$,
so $\chi_{\comm{A}{B}} - A\cdot\chi_B + B\cdot\chi_A\in\ssZ^{2,1}(\fa_-;\fa)$. But this space is zero by Corollary~\ref{coro:21-22-injective}. Hence $\delta$ satisfies \eqref{eq:delta-coycle-cond}. Thus we have shown the following.

\begin{Lemma}\label{lemma:invt-spencer-cocycle}
	Let $\mu\in\ssZ^{2,2}(\fa;\fa)$. Then $\mu|_{\fa_-\otimes\fa_-}\in\ssZ^{2,2}(\fa_-;\fa)$ with $\fh$-invariant cohomology class. If $V'=V$, we conversely have that if $\mu_-\in\ssZ^{2,2}(\fa_-;\fa)$ is a Spencer cocycle such that with $\fh$-invariant cohomology class then there exists a unique cocycle $\delta\in\ssZ^1(\fh;V^*\otimes\fh)\subseteq\Hom(\fh\otimes V,\fh)$ such that $\mu=\mu_-+\delta\in\ssZ^{2,2}(\fa;\fa)$.
\end{Lemma}

We can be more explicit about the form of $\delta$ in the $V=V'$ case: by Corollary~\ref{coro:21-22-injective}, there exists a unique map $\lambda\colon V\to\fso(V)$ such that
$\alpha(v,w) = \lambda(v)w - \lambda(w)v$.
One can then easily check that
$\delta(A,v) = \comm{A}{\lambda(v)} - \lambda(Av)
$
solves equations \eqref{eq:alpha-delta}--\eqref{eq:gamma-delta}.

\subsubsection{Jacobi identities}

We now consider the conditions imposed on the deforming maps by the Jacobi identities. These split into three types: those which are independent of the deforming maps, which are identically satisfied; those which involve only the deforming maps of degree~2, which are linear in those maps and which are equivalent to the homological conditions~\eqref{eq:22-spencer-cocycle-vss}--\eqref{eq:delta-coycle-cond}; and those which involve the degree-4 map $\theta$, all of which are quadratic in the deforming maps. We will not explicitly write the first two types, only the third type, since the equations are new. We denote by $[\fatilde_i,\fatilde_j,\fatilde_k]$ the Jacobi identity on the $(i, j,k)$-th homogeneous subspace of $\fa$. The Jacobi identities break down as follows:
\begin{itemize}\itemsep=0pt
	\item $[\fatilde_0,\fatilde_0,\fatilde_0]$: Identically satisfied by the Jacobi identity on $\fh\subseteq \fso(V)$.
	\item $[\fatilde_0,\fatilde_0,\fatilde_{-1}]$: Identically satisfied since $S'$ is a representation of $\fh$.
	\item $[\fatilde_0,\fatilde_0,\fatilde_{-2}]$: Equivalent to the cocycle condition \eqref{eq:delta-coycle-cond}, using that $V'$ is a representation of $\fh$.
	\item $[\fatilde_0,\fatilde_{-1},\fatilde_{-1}]$: Equivalent to cocycle condition \eqref{eq:gamma-delta}, using the $\fso(V)$-equivariance of $\kappa$.
	\item $[\fatilde_0,\fatilde_{-1},\fatilde_{-2}]$: Equivalent to cocycle condition~\eqref{eq:beta-delta}.
	\item $[\fatilde_0,\fatilde_{-2},\fatilde_{-2}]$: This identity has a $V'$-component, which is equivalent to the cocycle condition \eqref{eq:alpha-delta}, and an $\fh$-component which is quadratic in the deformation maps
		\begin{equation}\label{eq:jacobi-022}
			(A\cdot\theta)(v,w) = \delta(\delta(A,v),w) - \delta(\delta(A,w),v) - \delta(A,\alpha(v,w))
		\end{equation}
	for $A\in\fh$ and $v,w\in V'$.
	\item $[\fatilde_{-1},\fatilde_{-1},\fatilde_{-1}]$: Equivalent to the second Spencer cocycle condition \eqref{eq:22-spencer-cocycle-sss}.
	\item $[\fatilde_{-1},\fatilde_{-1},\fatilde_{-2}]$: Like the $[\fatilde_0,\fatilde_{-2},\fatilde_{-2}]$ identity, there is a $V'$-component and a $\fh$-component to this Jacobi identity. The former is the first Spencer cocycle condition \eqref{eq:22-spencer-cocycle-vss}, while the latter is the quadratic condition
		\begin{equation}\label{eq:jacobi-112}
			\theta(\kappa_s,v) + \delta(\gamma(s,s),v) + 2\gamma(s,\beta(v,s)) = 0
		\end{equation}
	for $s\in S'$ and $v\in V'$.
	\item $[\fatilde_{-1},\fatilde_{-2},\fatilde_{-2}]$: This identity gives the quadratic condition
		\begin{equation}\label{eq:jacobi-122}
			\theta(v,w)\cdot s + \beta(\alpha(v,w),s) - \beta(v,\beta(w,s)) + \beta(w,\beta(v,s)) =0,
		\end{equation}
	for $s\in S'$, $v,w\in V'$.
	\item $[\fatilde_{-2},\fatilde_{-2},\fatilde_{-2}]$: This has components in $V'$ and $\fh$, both of which are quadratic
		\begin{gather}
			\alpha(\alpha(u,v),w)
				+ \theta(u,v)w
				+ \text{cyclic perms.} = 0,\label{eq:jacobi-222a}\\
			\theta(\alpha(u,v),w)
				+ \delta(\theta(u,v),w)
				+ \text{cyclic perms.} =0\label{eq:jacobi-222b}
		\end{gather}
	for $u,v,w\in V'$.
\end{itemize}

In particular, recalling Lemma~\ref{lemma:invt-spencer-cocycle}, we have the following.

\begin{Proposition}\label{prop:gen-filtered-defs-cocycle}
	Let $\fa=V'\oplus S'\oplus\fh\subseteq\fs$ be a graded subalgebra, let $\alpha+\beta+\gamma+\delta\in\ssZ^{2,2}(\fa;\fa)$ and suppose we have a map $\theta\colon\Wedge^2V'\to\fh$. Then the brackets \eqref{eq:def-brackets-general} define a filtered deformation~$\fatilde$ of $\fa$ if and only if $\theta$ satisfies the equations~\eqref{eq:jacobi-022}--\eqref{eq:jacobi-222b}.
\end{Proposition}

Thus the task of determining the possible filtered deformations of a graded subalgebra $\fa\subseteq\fs$ consists of solving the linear equations \eqref{eq:22-spencer-cocycle-vss}--\eqref{eq:delta-coycle-cond} for $\mu=\alpha+\beta+\gamma+\delta$ and then the quadratic equations \eqref{eq:jacobi-022}--\eqref{eq:jacobi-222b} for $\theta$. We can interpret a $(2,2)$-cocycle $\mu$ as an \emph{infinitesimal filtered deformation} of $\fa$, and this deformation \emph{integrates} to an actual filtered deformation $\fatilde$ of $\fa$ if and only if there exists a solution $\theta$ to the equations \eqref{eq:jacobi-022}--\eqref{eq:jacobi-222b}. These equations are highly non-trivial in general, and even if an infinitesimal deformation integrates, the full deformation need not be unique. It is also difficult to give a homological characterisation of when two pairs~$(\mu,\theta)$ and $(\mu',\theta')$ give rise to isomorphic deformations.

\subsection{Homogeneity and high supersymmetry}
\label{sec:homogeneity}

We can get much better homological control over the problem of determining filtered deformations by passing to a particular class of graded subalgebras $\fa$. We have already seen in Lemmas~\ref{lemma:V-h-maps} and \ref{lemma:invt-spencer-cocycle} that taking $V=V'$ gives us significant simplifications. We would now like to guarantee that $S'$ is sufficiently large so that $V=\comm{S'}{S'}$, which will allow us to fully characterise filtered deformations of $\fa$ in terms of Spencer cohomology. To do this, we will make additional assumptions which allow us to use the homogeneity theorem, but we note that the essential ingredient for the results of this section is the fact that $V=\comm{S'}{S'}$.

\subsubsection{The homogeneity theorem}

We say that a symmetric Dirac current $\kappa$ is \emph{causal} if $\kappa_s=\kappa(s,s)$ is either timelike or null for all~${s\in S}$. In Lorentzian signature, a symmetric, causal Dirac current on the real pinor module $\ssS$ always exists in any dimension and can be used to construct Dirac currents on extended spinor modules. The existence of symmetric currents can be deduced from the Alekseevsky--Cort\'es classification \cite{Alekseevsky1997}; that they can be chosen to be causal must be shown case-by-case \cite{Hustler2016}. These can be used to build symmetric Dirac currents on ($N$-extended) spinor modules $S$ by restriction and tensoring with bilinears on an auxiliary module. However such a map is constructed, we have the following.

\begin{Theorem}[homogeneity theorem \cite{Figueroa-OFarrill2012,Hustler2016}]\label{thm:homogeneity}
	Let $(V,\eta)$ be a Lorentzian vector space of dimension $\dim V>2$ and $S$ a $($possibly $N$-extended$)$ spinor module of $\fso(V)$ with a symmetric, causal Dirac current $\kappa\colon \Odot^2 S \to V$. If $S'$ is a vector subspace of $S$ with $\dim S'>\frac{1}{2}\dim S$, then~${\kappa|_{\Odot^2S'}}$ is surjective onto $V$.
\end{Theorem}

For all $A\in\fso(V)$ and $s\in S$, $A\cdot s=0$ implies that $A \kappa_s=0$ by $\fso(V)$-equivariance of the Dirac current. Since $\fso(V)$ acts faithfully on $V$, the homogeneity theorem has the following corollary.

\begin{Corollary}\label{coro:homog-faithful}
	If the Dirac current $\kappa$ is causal and $\dim S'>\frac{1}{2}\dim S$, then the annihilator of~$S'$ in $\fso(V)$ is trivial. In particular, any subalgebra $\fh$ of $\fso(V)$ which preserves $S'$ acts faithfully on~$S'$.
\end{Corollary}

Guided by the homogeneity theorem (and by physics), let us now fix $\fs$ to be the Poincar\'e superalgebra associated to a Lorentzian inner product space $(V,\eta)$ and spinor module $S$ with symmetric causal Dirac current $\kappa$. Note that for such algebras we have $\comm{S'}{S'}=V$ for $\dim S'>\frac{1}{2}\dim S$, which prompts the following definition.

\begin{Definition}
	A graded subalgebra $\fa=V\oplus S'\oplus\fh$ of $\fs$ with $\dim S'> \frac{1}{2}\dim S$ is said to be \emph{highly supersymmetric}, and it is \emph{maximally supersymmetric} if $S'=S$.
	
	A filtered subdeformation $\fatilde$ of $\fs$ is said to be \emph{highly supersymmetric} if $\dim \fatilde_{\overline 1}> \frac{1}{2}\dim S$, and it is \emph{maximally supersymmetric} if $\dim \fatilde_{\overline 1} = \dim S$.
\end{Definition}

Clearly a filtered subdeformation $\fatilde$ is highly (maximally) supersymmetric if and only if its associated graded superalgebra $\fa$ is highly (maximally) supersymmetric as a graded subalgebra of $\fs$.

\begin{Remark}\label{rem:relaxing-homogeneity-assumption}
	The significance of the homogeneity theorem and the assumption of high supersymmetry is that they ensure that $\comm{S'}{S'}=V$, or $\comm{\fa_{-1}}{\fa_{-1}}=\fs_{-2}$. If one wishes to relax the assumptions which were made on the signature and Dirac current in order to invoke the theorem, one can simply impose this restriction on the graded subalgebra $\fa$ instead of the assumption of high supersymmetry and recover the results of the remainder of Sections~\ref{sec:filtered-def-poincare} and~\ref{sec:high-susy-spin-mfld}. One might term this \emph{sufficiently supersymmetry}.
\end{Remark}

\subsubsection{Homological properties}

Recall Definition~\ref{def:fund-trans-prolong} of some homological properties of graded Lie superalgebras. We now show that highly supersymmetric graded subalgebras of $\fs$ satisfy some of these.

\begin{Lemma}[\cite{Figueroa-OFarrill2017_1}]\label{lemma:high-susy-subalg-coho}
	Let $\fa$ be a highly supersymmetric graded subalgebra of $\fs$. Then
	\begin{itemize}\itemsep=0pt
	\item $\fa$ is fundamental and transitive,
	\item $\fa$ is a full prolongation of degree $2$,
	\item $\ssH^{d,2}(\fa_-;\fa)=0$ for all even $d>2$.
	\end{itemize}
\end{Lemma}
\begin{proof}
	Theorem~\ref{thm:homogeneity} precisely says that $\fa$ is fundamental, and transitivity follows from faithfulness of the action of $\fh=\fa_0\subseteq\fso(V)$ on $V=\fa_{-2}$ or $S'=\fa_{-1}$. Corollary~\ref{coro:21-22-injective} then says that~${\ssH^{2,1}(\fa_-;\fa)=0}$. For $d>2$, $\ssH^{d,1}(\fa_-;\fa)=0$ trivially since $\ssC^{d,1}(\fa_-;\fa)=0$. The beginning of the degree-4 Spencer complex is
	\begin{align*}%\label{eq:spencer-4}
		0	&\longrightarrow \Hom\bigl(\Wedge^2V,\fh\bigr)	
			\longrightarrow \Hom\bigl(\Wedge^3 V,V\bigr)\oplus \Hom\bigl(\Wedge^2V\otimes S',S'\bigr) \oplus \Hom\bigl(V\otimes \Odot^2 S', \fh\bigr) \\
&\longrightarrow \cdots,
	\end{align*}
	and for $\theta\in \ssC^{4,2}(\fa_-;\fa)=\Hom\bigl(\Wedge^2V,\fh\bigr)$, we have
$\partial\theta (v,s,s) = \theta(v,\kappa_s)$
	for all $v\in V$, $s\in S'$, so again by homogeneity, $\partial\theta=0$ if and only if $\theta=0$. This shows that $\ssH^{d,4}(\fa_-;\fa)=0$, and we have~${\ssH^{d,2}(\fa_-;\fa)=0}$ for all $d>4$ since $\ssC^{d,2}(\fa_-;\fa)=0$.
\end{proof}

The following pivotal observation is a straightforward generalisation to all spacetime dimensions of \cite[Proposition 6]{Figueroa-OFarrill2017_1} and also of the analogous results for \emph{maximally} supersymmetric graded subalgebras (those for which $\fa_{-1}=S$) in \cite{deMedeiros2016,Figueroa-OFarrill2017}. For the proof, we recall Propositions~\ref{prop:def-seq-spencer-cocycle} and \ref{prop:def-seq-redef} which characterise the defining sequences of filtered deformations in terms of Spencer cohomology and tell us how much freedom we have to redefine them.

\begin{Proposition}[\cite{Figueroa-OFarrill2017_1}]\label{prop:highly-susy-filtered-def}
	Let $\fa=V'\oplus S'\oplus\fh$ be a graded subalgebra of $\fs$ and $\fatilde$ a filtered deformation of $\fa$ with defining sequence $(\mu,\theta,0,\dots)$ as above. Then
	\begin{enumerate}\itemsep=0pt
		\item[$(1)$] $\mu|_{\fa_-\otimes\fa_-}$ is a cocycle in $\ssC^{2,2}(\fa_-;\fa)$, and
$[\mu|_{\fa_-\otimes\fa_-}] \in \ssH^{2,2}(\fa_-;\fa)^\fh$.
		Furthermore, $\mu$ is a~cocycle in $\ssC^2(\fa;\fa)$.	%\label{item:filtered-subdef-cocycle}
		\item[$(2)$] If $\fatilde$ is highly supersymmetric and $\fatilde'$ is another filtered deformation of $\fa$ with degree-$2$ deformation map $\mu'$ such that $[\mu'|_{\fa_-\otimes\fa_-}]=[\mu|_{\fa_-\otimes\fa_-}]$ then $\fatilde\cong\fatilde'$ as filtered Lie superalgebras.	%\label{item:high-susy-filtered-subdef-cocycle}
	\end{enumerate}
\end{Proposition}

\begin{proof}
	The first statement follows by Proposition~\ref{prop:def-seq-spencer-cocycle}. For the second, if $[\mu'|_{\fa_-\otimes\fa_-}]=[\mu|_{\fa_-\otimes\fa_-}]$ then $(\mu'-\mu)|_{\fa_-\otimes\fa_-}$ is a Spencer coboundary, so by point \ref{item:def-seq-redef-1} of Proposition~\ref{prop:def-seq-redef} there is a new defining sequence $\{\mu'',\theta'',0,\dots\}$ of $\fatilde'$ such that $\mu''|_{\fa_-\otimes\fa_-}=\mu|_{\fa_-\otimes\fa_-}$. Then, since $\fa$ is a full prolongation of degree 2 by Lemma~\ref{lemma:high-susy-subalg-coho}, by point \ref{item:def-seq-redef-3} of Proposition~\ref{prop:def-seq-redef} we can again find a new defining sequence $\{\mu''',\theta''',0,\dots\}$ of $\fatilde'$ such that $\mu'''=\mu$ and $\theta'''=\theta'''|_{\fa_-\otimes\fa_-}=\theta|_{\fa_-\otimes\fa_-}=\theta$.
\end{proof}

Note that we already saw the first statement of the proposition above in Section~\ref{sec:filtered-def-cohom}. The second statement allows us to strengthen Proposition~\ref{prop:gen-filtered-defs-cocycle}; by Lemma~\ref{lemma:invt-spencer-cocycle}, a choice of $\fh$-invari\-ant~$\mu_{\fa_-\otimes\fa_-}$ determines $\mu$ for $V=V'$; we now see that the cohomology class $[\mu_{\fa_-\otimes\fa_-}]$ determines the entire deformation up to isomorphism in the highly supersymmetric case.

What this means in practice is that in order to find the infinitesimal deformations of a~highly supersymmetric graded subalgebra $\fa$, we compute $\ssH^{2,2}(\fa_-;\fa)$ and its $\fh$-invariants. The other deforming maps $\delta$ and $\theta$ will then be determined by the Jacobi identities. Note, though, that the Jacobi identities which are quadratic in the deformation maps may still obstruct the ``integration'' of the infinitesimal deformation $[\mu_{\fa_-\otimes\fa_-}]$ to a filtered deformation $\fatilde$ of $\fa$.

\subsubsection{Reparametrising the deformation}
\label{sec:reparam-defs}

It will be useful for us to relate $\ssH^{2,2}(\fa_-;\fa)$ to some other cohomology groups. The inclusion $i\colon\fa\hookrightarrow\fs$ of a graded subalgebra into $\fs$ induces the following maps of cochains
\begin{align*}
	i_*\colon\ \ssC^{\bullet,\bullet}(\fa_-;\fa) \to \ssC^{\bullet,\bullet}(\fa_-;\fs),
		\qquad \phi\mapsto i_*\phi = i\circ\phi,\\
	i^*\colon\ \ssC^{\bullet,\bullet}(\fs_-;\fs) \to \ssC^{\bullet,\bullet}(\fa_-;\fs),
		\qquad \phi\mapsto i^*\phi = \phi\circ i,
\end{align*}
where $\fs$ is viewed as representation of $\fa_-$ via the restriction of the adjoint representation, and the cochain complex $\ssC^{\bullet,\bullet}(\fa_-;\fs)$ has been given a grading in the same manner as the Spencer complex (equation~\eqref{eq:spencer-complex-grading}); the differential respects this grading and the maps $i_*$, $i^*$ respect both this grading and the homological grading -- that is, they commute with the differential $\partial$. Thus these maps in turn induce maps in cohomology also denoted $i_*$ and $i^*$.

Let us consider these maps in degree $(2,2)$ for the highly supersymmetric case. Recall that, by Lemma~\ref{lemma:H22-full-flat-model}, $\ssH^{2,2}(\fs_-;\fs)\cong\cH^{2,2}$ where the latter is the space of normalised Spencer $(2,2)$-cocycles of $\fs$.

\begin{Lemma}\label{lemma:maps-in-cohomology}
	Let $\fa$ be a highly supersymmetric graded subalgebra of $\fs$. Then we have a diagram with exact rows and columns
	\[
	\begin{tikzcd}
		& & & 0 \arrow[d]\\
		& & & \ssH^{2,2}(\fa_-;\fa) \arrow[d,"i_*"] \\
		0 \arrow[r]
		& \cK^{2,2}(\fa_-) \arrow[r]
		& \ssH^{2,2}(\fs_-;\fs) \cong \cH^{2,2} \arrow[r, "i^*"]
		& \ssH^{2,2}(\fa_-;\fs),
	\end{tikzcd}
	\]
	where
$\cK^{2,2}(\fa_-) := \big\{ \beta+\gamma\in \cH^{2,2} \mid \beta|_{V\otimes S'}=0 \big\} \cong \ker i^*$.
\end{Lemma}

\begin{proof}
	For $\beta\colon V\otimes S\to S$, $\gamma\colon \Odot^2S\to\fso(V)$ such that $\beta+\gamma\in\cH^{2,2}$, we have $i^*[\beta+\gamma]=0$ if and only if $\beta|_{V\otimes S'}+\gamma|_{\Odot^2S'}=\partial\lambda$, where $\lambda\colon V\to\fso(V)$. This implies that $\partial\lambda|_{\Wedge^2 V}=0$, so $\lambda=0$ by Lemma~\ref{lemma:V-h-maps}, thus $\beta|_{V\otimes S'}=0$ and $\gamma|_{\Odot^2S'}=0$. Note that $\beta|_{V\otimes S'}=0$ implies $\gamma|_{\Odot^2S'}=0$ by the cocycle condition for $\beta+\gamma$, so the isomorphism $\cH^{2,2}\cong\ssH^{2,2}(\fs_-;\fs)$ restricts to give $\cK^{2,2}(\fa_-)\cong \ker i^*$.
	
	For $\alpha\colon \Wedge^2V\to V$, $\beta\colon V\otimes S'\to S'$, $\gamma\colon \Odot^2S'\to\fh$ with $\alpha+\beta+\gamma\in\ssZ^{2,2}(\fa_-;\fa)$, $i_*[\alpha+\beta+\gamma]=0$ if and only if $i_*(\alpha+\beta+\gamma)=\partial\lambda$ where $\lambda:V\to\fso(V)$. Then $\gamma(s,s)=-\lambda(\kappa(s,s))$ for all $s\in S'$; in particular, $\lambda(V)=\lambda(\comm{S'}{S'})\subseteq\fh$, or equivalently $\lambda=i_*\lambda'$ for some $\lambda'\in\ssC^{2,1}(\fa_-;\fa)=\Hom(V,\fh)$. It follows that $\alpha+\beta+\gamma=\partial\lambda'$, so $[\alpha+\beta+\gamma]=0$, showing that $i_*\colon \ssH^{2,2}(\fa_-;\fa)\to \ssH^{2,2}(\fa_-;\fs)$ is injective.
\end{proof}

By the same argument as in the proof of Lemma~\ref{lemma:H22-full-flat-model}, in the highly supersymmetric case any cohomology class in $\ssH^{2,2}(\fa_-;\fs)$ has a unique representative in the space of normalised cocycles
\[%\label{eq:norm-cocycles-fa;fs}
	\cH^{2,2}(\fa_-)
	= \bigl\{\beta+\gamma \in \ssZ^{2,2}(\fa_-;\fs)
		\mid
			\beta\in\Hom(V\otimes S',S),\,
			\gamma\in\Hom\bigl(\Odot^2S',\fso(V)\bigr)\bigr\},
\]
and we have a splitting of $\fh$-modules $\ssZ^{2,2}(\fa_-;\fs)=\cH^{2,2}(\fa_-)\oplus\ssB^{2,2}(\fa_-;\fs)$, whence as $\fh$-modules $\ssH^{2,2}(\fa_-;\fs) \cong \cH^{2,2}(\fa_-)$. We thus have an exact sequence at the level of cocycles corresponding to the bottom row of the diagram in the lemma above
\[
\begin{tikzcd}
	0 \arrow[r]
	& \cK^{2,2}(\fa_-) \arrow[r]
	& \cH^{2,2} \arrow[r, "i^*"]
	& \cH^{2,2}(\fa_-).
\end{tikzcd}
\]
On the other hand, any cohomology class in $\ssH^{2,2}(\fa_-;\fa)$ defines a unique normalised cocycle in~$\cH^{2,2}(\fa_-)$ by pushing forward to $\ssH^{2,2}(\fa_-;\fs)$ and then identifying the normalised cocycle. Since~$i_*$ is injective in cohomology, if a cocycle in $\cH^{2,2}(\fa_-)$ is obtained from a class in $\ssH^{2,2}(\fa_-;\fa)$, that class is unique. We can be much more explicit about this: if $\alpha+\beta+\gamma\in\ssZ^{2,2}(\fa_-;\fa)$ then there exists a unique element $\lambda\in\ssC^{2,2}(\fa_-;\fs)=\Hom(V,\fso(V))$ such that
\begin{equation}\label{eq:alpha-lambda}
	\alpha(v,w) = \lambda(v)w - \lambda(w)v;
\end{equation}
the components $\betatilde\colon V\otimes S' \to S$ and $\gammatilde\colon \Odot^2 S' \to \fso(V)$ of the unique normalised cocycle $\betatilde+\gammatilde\in\cH^{2,2}(\fa_-)$ such that $i_*[\alpha+\beta+\gamma]=[\betatilde+\gammatilde]$ are then given by
\begin{align}%\label{eq:beta-tilde}
	\betatilde(v,s) = \beta(v,s) - \lambda(v)\cdot s,\qquad
	\gammatilde(s,s) = \gamma(s,s) + \lambda(\kappa_s), \label{eq:gamma-tilde}
\end{align}
for $v\in V$, $s\in S'$. The maps $\lambda$, $\betatilde$, $\gammatilde$ are unique for a given cocycle $\alpha+\beta+\gamma$, while a different choice of representative for the same cohomology class corresponds to shifting $\lambda$ by a map $V\to\fh$. We note that
$i_*(\alpha+\beta+\gamma) = \betatilde+\gammatilde+\partial\lambda$,
and that we have the constraints
$\betatilde(v,s) + \lambda(v)\cdot s \in S'$, $
	\gammatilde(s,s) - \lambda(\kappa_s) \in \fh$
for all $v\in V$, $s\in S'$. Finally, observe that $[\alpha+\beta+\gamma]$ is $\fh$-invariant if and only if $\betatilde+\gammatilde\in\cH^{2,2}(\fa_-)$ is (which is the case if and only if $\betatilde$ is $\fh$-invariant). If this is the case, the discussion following Lemma~\ref{lemma:invt-spencer-cocycle} tells us that the map $\delta\colon \fh\otimes V\to\fh$ defined by
\begin{equation}\label{eq:delta-lambda}
	\delta(A,v) = \comm{A}{\lambda(v)} - \lambda(Av)
\end{equation}
is the unique such map so that $\alpha+\beta+\gamma+\delta\in\ssZ^{2,2}(\fa;\fa)$.

The discussion above allows us to reparametrise the deformed brackets \eqref{eq:def-brackets-general} in the highly supersymmetric case.

\begin{Proposition}\label{prop:filtered-def-high-susy}
	If $\fatilde$ is a highly supersymmetric filtered deformation of a graded subalgebra~${\fa=V\oplus S'\oplus \fh}$ of $\fs$, then $\fatilde$ has a presentation of the form
	\begin{gather}
		\comm{A}{B} = \underbrace{AB-BA}_\fh,\qquad
		\comm{s}{s} = \underbrace{\kappa_s}_V + \underbrace{\gammatilde(s,s) - \lambda(\kappa_s)}_\fh,\nonumber\\
		\comm{A}{v} = \underbrace{Av}_V + \underbrace{\comm{A}{\lambda(v)} - \lambda(Av)}_\fh,\qquad
		\comm{v}{s} = \underbrace{\betatilde(v,s) + \lambda(v)\cdot s}_{S'},\qquad
		\comm{A}{s} = \underbrace{A\cdot s}_{S'},\nonumber
	\\
		\comm{v}{w} = \underbrace{\lambda(v)w-\lambda(w)v}_V + \underbrace{\thetatilde(v,w) - \lambda(\lambda(v)w-\lambda(w)v) + \comm{\lambda(v)}{\lambda(w)}}_\fh\label{eq:def-brackets-V}
	\end{gather}
	for $A,B\in\fh$, $v,w\in V$, $s\in S'$, where the deforming maps are
	\[
\betatilde\colon\ V\otimes S' \to S,\qquad
		\gammatilde\colon\ \Odot^2 S' \to \fso(V),\qquad
		\lambda\colon\ V\to\fso(V),\qquad
		\thetatilde\colon\ \Wedge^2 V \to \fso(V),
	\]
	where $\betatilde+\gammatilde\in\cH^{2,2}(\fa_-)^\fh$, $\thetatilde$ is $\fh$-invariant and the remaining Jacobi identities are equivalent to
	\begin{gather}
		\thetatilde(\kappa_s,v) + 2\gammatilde\bigl(s,\betatilde(v,s)+\lambda(v)\cdot s\bigr) - \comm{\lambda(v)}{\gammatilde(s,s)} = 0, \label{eq:tilde-theta-kappa}\\
		\thetatilde(v,w)\cdot s
		= \betatilde\bigl(v,\betatilde(w,s)+\lambda(w)\cdot s\bigr)
			- \lambda(w)\cdot \betatilde(v,s) + \betatilde(\lambda(w)v,s)\nonumber\\
			\phantom{\thetatilde(v,w)\cdot s
		=}{} - \betatilde\bigl(w,\betatilde(v,s)+\lambda(v)\cdot s\bigr)
			+ \lambda(v)\cdot \betatilde(w,s) - \betatilde(\lambda(v)w,s),\label{eq:tilde-theta-lambda}\\
		\thetatilde(u,v)w + \thetatilde(v,w)u + \thetatilde(w,u)v =0,\label{eq:tilde-theta-bianchi}\\
		\bigl(\lambda(u)\cdot\thetatilde\bigr)(v,w) + \bigl(\lambda(v)\cdot\thetatilde\bigr)(w,u) +			\bigl(\lambda(w)\cdot\thetatilde\bigr)(u,v) = 0 \label{eq:tilde-theta-lambda-bianchi}
	\end{gather}
	for all $u,v,w\in V$ and $s\in S'$.
	Two deformations $\fatilde$, $\fatilde'$ described in this way are isomorphic as filtered Lie superalgebras if and only if they determine the same element $\betatilde+\gammatilde\in\cH^{2,2}(\fa_-)^\fh$.
\end{Proposition}

\begin{proof}
	We begin with a presentation of the form \eqref{eq:def-brackets-general}, so that $\alpha+\beta+\gamma$ is a $(2,2)$-Spencer cocycle for $\fa$ with $\fh$-invariant cohomology class and then define $\lambda,\betatilde,\gammatilde$ by equations \eqref{eq:alpha-lambda}--\eqref{eq:gamma-tilde}; $\delta$ must then be given by \eqref{eq:delta-lambda}. The brackets then take the claimed form if we define
$\thetatilde(v,w) = \theta(v,w) + \lambda(\lambda(v)w-\lambda(w)v) - \comm{\lambda(v)}{\lambda(w)}$.
	It remains to check the Jacobi identities~\eqref{eq:jacobi-022}--\eqref{eq:jacobi-222b} in our new variables. Substituting the definition of $\delta$ into the first identity \eqref{eq:jacobi-022} gives us $A\cdot \thetatilde=0$ for all $A\in\fh$. Next, we find that~\eqref{eq:jacobi-112} implies \eqref{eq:tilde-theta-kappa}
	\begin{gather*}
		\theta(\kappa_s,v) + \delta(\gamma(s,s),v) + 2\gamma(s,\beta(v,s)) \\
		\qquad= \theta(\kappa_s,v) + \comm{\gammatilde(s,s)}{\lambda(v)} - \lambda(\gammatilde(s,s)v) - \comm{\lambda(\kappa_s)}{\lambda(v)} + \lambda(\lambda(\kappa_s)v) \\
		\phantom{\qquad=}{}+ 2\gammatilde(s,\betatilde(v,s)+\lambda(v)\cdot s) - 2\lambda(\kappa(s,\betatilde(v,s)+\lambda(v)\cdot s)) \\
	\qquad	= \theta(\kappa_s,v) + \lambda\qty(\lambda(\kappa_s)v - 2\kappa(s,\lambda(v)\cdot s)) - \comm{\lambda(\kappa_s)}{\lambda(v)} + 2\gammatilde(s,\betatilde(v,s)+\lambda(v)\cdot s)\\
			\phantom{\qquad=}{}- \comm{\lambda(v)}{\gammatilde(s,s)}
			- \lambda(\kappa(s,s)v + 2\kappa(s,\betatilde(v,s)))\\
		\qquad= \thetatilde(\kappa_s,v) + 2\gammatilde(s,\betatilde(v,s)+\lambda(v)\cdot s) - \comm{\lambda(v)}{\gammatilde(s,s)},
	\end{gather*}
	where in the last line we have used the definition of $\thetatilde$, a cocycle condition for $\betatilde+\gammatilde$ and $\fso(V)$-invariance of $\kappa$. Substitution of the change of variables followed by some trivial rearrangement shows that \eqref{eq:jacobi-122} and \eqref{eq:jacobi-222a} are equivalent to \eqref{eq:tilde-theta-lambda} and \eqref{eq:tilde-theta-bianchi} respectively. Finally, \eqref{eq:jacobi-222b} gives~\smash{$\bigl(\lambda(u)\cdot\thetatilde\bigr)(v,w)
			+ \lambda\bigl(\thetatilde(u,v)w\bigr)
	+ \text{cyclic perms.}= 0$},
	but then substitution of \eqref{eq:tilde-theta-bianchi} yields~\eqref{eq:tilde-theta-lambda-bianchi}. The final claim follows by Proposition~\ref{prop:highly-susy-filtered-def} since the cocycle $\betatilde+\gammatilde\in\cH^{2,2}(\fa_-)^\fh$ is uniquely determined by the cohomology class $[\alpha+\beta+\gamma]\in\ssH^{2,2}(\fa_-;\fa)$.
\end{proof}

\subsection{Realisable subdeformations}
\label{sec:realisable-subdef}

We have just seen that if $\alpha+\beta+\gamma\in\ssZ^{2,2}(\fa_-;\fa)$ is the Spencer cocycle associated to some presentation of a filtered deformation of a highly supersymmetric graded subalgebra $\fa\subseteq\fs$, then there is an $\fh$-invariant cocycle $\betatilde+\gammatilde\in\ssZ^{2,2}(\fa_-;\fs)$ such that $i_*[\alpha+\beta+\gamma]=[\betatilde+\gammatilde]$. We will now restrict our attention to cocycles obeying a stronger version of this and the corresponding filtered subdeformations of $\fs$ (where they exist). Essentially, we demand that $\betatilde+\gammatilde$ be the pull-back of some $\fh$-invariant element $\betahat+\gammahat\in\cH^{2,2}$.

\subsubsection{Admissibility}

\begin{Definition}\label{def:admissible-class}
	A cohomology class $[\alpha+\beta+\gamma]\in \ssH^{2,2}(\fa_-;\fa)$ for a highly supersymmetric graded subalgebra $\fa$ of $\fs$ is \emph{admissible} if we have
$i_*[\alpha+\beta+\gamma]=i^*\big[\betahat+\gammahat\big]
$
	in $\ssH^{2,2}(\fa_-;\fs)$ for some $\fh$-invariant cocycle \smash{$\betahat+\gammahat\in \cH^{2,2}$}.
	
	A cocycle in $\ssZ^{2,2}(\fa_-;\fa)$ is \emph{admissible} if its cohomology class is admissible.
\end{Definition}

Note that an admissible cohomology class is in particular $\fh$-invariant. Unpacking the definition a little, for admissible cocycles we have
\begin{equation}\label{eq:admiss-cocycle-param}
	i_*(\alpha + \beta + \gamma) = i^*(\betahat + \gammahat) + \partial\lambda
\end{equation}
for a unique $\lambda\in\ssZ^{2,2}(\fa_-;\fs)=\Hom(V,\fso(V))$ and $\betahat+\gammahat\in\bigl(\cH^{2,2}\bigr)^\fh$ determined uniquely up to elements in $\cK^{2,2}(\fa_-)$. If we fix an admissible cohomology class but not a particular representative, $\lambda$ is only determined up to elements in $\ssZ^{2,2}(\fa_-;\fa)=\Hom(V,\fh)$.

Recalling that a highly supersymmetric subdeformation of $\fs$ uniquely determines an element of $\ssH^{2,2}(\fa_-;\fa)^\fh$ by Proposition~\ref{prop:highly-susy-filtered-def}, we make the following definition.

\begin{Definition}[geometric realisability]\label{def:geom-real}
	A highly supersymmetric filtered subdeformation $\fatilde$ of $\fs$ is \emph{$($geometrically$)$ realisable} if the associated cohomology class in $\ssH^{2,2}(\fa_-;\fa)^\fh$ is admissible.
\end{Definition}

The definitions above are inspired by, although \emph{not} equivalent to \cite[Definition~9]{Figueroa-OFarrill2017_1} in the context of 11-dimensional supergravity. There, one finds that $\cH^{2,2}\cong \Wedge^4V$, and it is the 4-form~$\varphi$ associated to $\betahat+\gammahat$ which is called admissible if the conditions above are met \emph{and} an additional algebraic condition related to the fact that $\varphi$ should considered to be the value at a point of a~closed 4-form field strength on some supergravity background is imposed. Our definition here is weaker because it is not clear \textit{a priori} how this condition should be generalised.

Let us make some particular remarks about these definitions in the maximally supersymmetric case, where $\fa_-=\fs_-$. Here, the map $i^*\colon\ssC^{2,2}(\fs_-;\fs)\to\ssC^{2,2}(\fa_-;\fs)$ whose induced map in cohomology appears in the definition of admissibility, is the identity map. An admissible cocycle for $\fa$ is then simply a cocycle in $\ssZ^{2,2}(\fa_-;\fa)$ whose cohomology class is $\fa_0$-invariant. It follows that \emph{all} maximally supersymmetric filtered deformations of $\fs$ are realisable.

We will justify the term ``geometrically realisable'' in Section~\ref{sec:high-susy-spin-mfld-homogeneous} (see Theorem~\ref{thm:homog-spin-mfld}). Our goal in this section will be to determine the conditions under which an admissible cohomology class of a highly supersymmetric graded subalgebra $\fa\subseteq\fs$ is \emph{integrable}, in the sense that there exists a corresponding geometrically realisable filtered deformation $\fatilde$. We first record the following. Note that the homogeneity theorem (see Theorem~\ref{thm:homogeneity}) is required for this result as well as for many in Section~\ref{sec:integrability}.

\begin{Lemma}\label{lemma:admiss-props}
	Let $\alpha+\beta+\gamma$ be an admissible cocycle parametrised as in \eqref{eq:admiss-cocycle-param}. Then
	\begin{itemize}\itemsep=0pt
		\item for all $v\in V$, $s\in S'$, $\betahat(v,s)+\lambda(v)\cdot s \in S'$;
		\item for all $s,s'\in S'$, $\gammahat(s,s')-\lambda(\kappa(s,s'))\in\fh$;
		\item for all $A\in\fh$, $v\in V$, $\comm{A}{\lambda(v)}-\lambda(Av)\in\fh$.
	\end{itemize}
\end{Lemma}

\begin{proof}
	The first two claims follow directly from the definition. For the third, by Theorem~\ref{thm:homogeneity} we can take $v=\kappa_s$ for some $s\in S'$. Recalling that the action of $\fh$ preserves $S'$ and using~$\fso(V)$-equivariance of $\kappa$ and $\fh$-invariance of $\gammahat$, we have
	\begin{align*}
		\comm{A}{\lambda(\kappa_s)} - \lambda(A\kappa_s)
		 &= \comm{A}{\lambda(\kappa_s)} - 2\lambda(\kappa(As,s))\\
		 &= \comm{A}{\lambda(\kappa_s)-\gamma(s,s)} + 2(\gammahat(As,s)-\lambda(\kappa(As,s))),
	\end{align*}
	and both of the terms in the last expression lie in $\fh$ by the second claim of the lemma.
\end{proof}

\subsubsection{Integrability}
\label{sec:integrability}

Suppose we have an admissible cocycle parametrised as in \eqref{eq:admiss-cocycle-param} and take $\betatilde=i^*\betahat$, $\gammatilde=i^*\gammahat$. We wish to know whether it is possible to define a (geometrically realisable) filtered subdeformation as in Proposition~\ref{prop:filtered-def-high-susy} using these maps in the definition of the bracket \eqref{eq:def-brackets-V}. By the above lemma, it remains only to check whether there exists an $\fh$-invariant map $\thetatilde\colon \Wedge^2V\to \fso(V)$ such that
\begin{equation}\label{eq:theta-lambda-in-h}
	\thetatilde(v,w) - \lambda(\lambda(v)w-\lambda(w)v) + \comm{\lambda(v)}{\lambda(w)} \in \fh
\end{equation}
satisfying the equations \eqref{eq:tilde-theta-kappa} to \eqref{eq:tilde-theta-lambda-bianchi}. We first observe that equations~\eqref{eq:tilde-theta-kappa} and \eqref{eq:tilde-theta-lambda} simplify slightly
\begin{gather}
	\thetatilde(v,\kappa_s) = 2\gammahat\bigl(s,\betahat(v,s)\bigr) - \bigl(\lambda(v)\cdot\gammahat\bigr)(s,s),
	\label{eq:tilde-theta-def-1}\\
		\thetatilde(v,w)\cdot s
		= \betahat\bigl(v,\betahat(w,s)\bigr) - \betahat\bigl(w,\betahat(v,s)\bigr) + \bigl(\lambda(v)\cdot\betahat\bigr)(w,s) - \bigl(\lambda(w)\cdot\betahat\bigr)(v,s),
		\label{eq:tilde-theta-act-s}
\end{gather}
and by Theorem~\ref{thm:homogeneity}, either equation fully determines $\thetatilde$ (if it exists). In fact, $\thetatilde$ is uniquely determined by the admissible cohomology class: it depends only on $i^*\betahat$, $i^*\gammahat$ and $\lambda$, so a shift in~${\betahat+\gammahat}$ by an element of $\cK^{2,2}(\fa_-)$ does not change it, while by $\fh$-invariance of $\betahat$ and $\gammahat$, changing~$\lambda$ by addition of a map $V\to\fh$ also does not change $\thetatilde$. There may, however, be obstructions to the existence of $\thetatilde\colon\Wedge^2 V\to \fso(V)$ even for an admissible cocycle. For now, we will postpone considering this issue and show that a map $\thetatilde$ determined by the equations above automatically satisfies all of the other required properties.

\begin{Lemma}\label{lemma:theta-lambda-in-h}
	Any map $\thetatilde\colon\Wedge^2V\to \fso(V)$ satisfying equation~\eqref{eq:tilde-theta-def-1} also satisfies~\eqref{eq:theta-lambda-in-h}.
\end{Lemma}

\begin{proof}
	By Theorem~\ref{thm:homogeneity}, it is sufficient to show the result for $w=\kappa_s$ for some $s\in S'$; we have
	\begin{gather*}
		\thetatilde(v,\kappa_s) - \lambda(\lambda(v)\kappa_s-\lambda(\kappa_s)v) + \comm{\lambda(v)}{\lambda(\kappa_s)}\\
			\qquad= 2\big\{\qty(\gammahat-\lambda\circ\kappa)(s,\betahat(v,s)+\lambda(v)\cdot s)\big\}\\
				\phantom{\qquad= }{} + \qty{\comm{(\gammahat-\lambda\circ\kappa)(s,s)}{\lambda(v)} - \lambda((\gammahat-\lambda\circ\kappa)(s,s)v)},
	\end{gather*}
	and since $(\gammahat-\lambda\circ\kappa)|_{\Odot^2S'}$ takes values in $\fh$, Lemma~\ref{lemma:admiss-props} implies that both of the expressions within the braces lie in $\fh$.
\end{proof}

\begin{Lemma}\label{lemma:tilde-theta-def-2}
	Any map $\thetatilde\colon\Wedge^2 V\to \fso(V)$ satisfying equation~\eqref{eq:tilde-theta-act-s} also satisfies
	\begin{gather*}
		\thetatilde(v,w)\kappa_s
			= 2\gammahat\bigl(\betahat(v,s),s\bigr)w
				- 2\gammahat\bigl(\betahat(w,s),s\bigr)v- \bigl(\lambda(v)\cdot\gammahat\bigr)(s,s)w
				+ \qty(\lambda(w)\cdot\gammahat)(s,s)v\label{eq:tilde-theta-def-2}
	\end{gather*}
	for all $v,w\in V$, $s\in S'$, and this equation determines $\thetatilde$ uniquely.
\end{Lemma}

\begin{proof}
	First, it will be useful to depolarise one of the cocycle conditions for $\betahat+\gammahat$
	\begin{equation}\label{eq:22-spencer-cocycle-vss-depol}
		2\kappa\bigl(\betahat(v,s),s'\bigr) + 2\kappa\bigl(\betahat(v,s'),s\bigr) + 2\gammahat(s,s')v = 0.
	\end{equation}
	Using $\fso(V)$-invariance of $\kappa$, equation~\eqref{eq:tilde-theta-act-s} gives
	\begin{align*}
		\thetatilde(v,w)\kappa_s
		= 2\kappa\bigl(\thetatilde(v,w)\cdot s,s\bigr)
		={}& 2\kappa\bigl(\betahat(v,\betahat(w,s)\bigr),s)
			+ 2\kappa\bigl(\lambda(v)\cdot\betahat(w,s),s\bigr)\\
			& - 2\kappa\bigl(\betahat(\lambda(v)w,s),s\bigr)
			- 2\kappa\bigl(\betahat(w,\lambda(v)\cdot s),s\bigr)
			- (v \leftrightarrow w).
	\end{align*}
	We now use equation~\eqref{eq:22-spencer-cocycle-vss-depol} with $s'=s,\lambda(v)\cdot s$ and $\betahat(w,s) $ to replace some of the terms above and then use $\fso(V)$-invariance of $\kappa$ again to conclude that
	\begin{align*}
		\thetatilde(v,w)\kappa_s
		={}& - 2\gammahat\bigl(\betahat(w,s),s\bigr)v
			- 2\kappa\bigl(\betahat(v,s),\betahat(w,s)\bigr) + 2\kappa\bigl(\lambda(v)\cdot\betahat(w,s),s\bigr)
\\
			&+ \gammahat(\lambda(v)\cdot s,s)v + \gammahat(s,s)(\lambda(v)w)
			+ 2\kappa\bigl(\betahat(w,s),\lambda(v)\cdot s\bigr)
			- (v \leftrightarrow w)\\
		={}&2\gammahat\bigl(\betahat(v,s),s\bigr)w
			+ \lambda(v)\kappa\bigl(\betahat(w,s),s\bigr)\\
			& + \gamma(s,s)(\lambda(v)w)
			+ \gammahat(\lambda(v)\cdot s,s)w - (v \leftrightarrow w),
	\end{align*}
	which is the desired expression. In the final equality, we have used the skew-symmetry of the expression in $v$ and $w$. By Theorem~\ref{thm:homogeneity}, this defines $\thetatilde$ uniquely.
\end{proof}

\begin{Lemma}\label{lemma:quadr-jacobi-dependency}
	For any map $\thetatilde\colon\Wedge^2 V\to\fso(V)$, we have the following:
	\begin{itemize}\itemsep=0pt
		\item Either equation~\eqref{eq:tilde-theta-def-1} or \eqref{eq:tilde-theta-act-s} implies that $\thetatilde$ is $\fh$-invariant.
		\item Assuming equation~\eqref{eq:tilde-theta-def-1}, equation~\eqref{eq:tilde-theta-def-2} is equivalent to the algebraic Bianchi identity~\eqref{eq:tilde-theta-bianchi}. In particular, equations \eqref{eq:tilde-theta-def-1} and \eqref{eq:tilde-theta-act-s} together imply the Bianchi identity.
		\item Equations~\eqref{eq:tilde-theta-def-1} and \eqref{eq:tilde-theta-act-s} together imply \eqref{eq:tilde-theta-lambda-bianchi}.
	\end{itemize}
\end{Lemma}

\begin{proof}
	The first claim follows by some fairly straightforward but long manipulations which use $\fso(V)$-invariance of $\kappa$, $\fh$-invariance of $\betahat$ and $\gammahat$ and the fact that $\lambda(Av)-\comm{A}{\lambda(v)}\in\fh$ for all~${A\in\fh}$. For the second claim, let us assume equation \eqref{eq:tilde-theta-def-1} holds. Then equation~\eqref{eq:tilde-theta-def-2} gives~${\thetatilde(v,w)\kappa_s = \thetatilde(v,\kappa_s)w - \thetatilde(w,\kappa_s)v}
$
	for all $v,w\in V$, $s\in S'$, which by the homogeneity theorem is the Bianchi identity. Conversely, if we assume the Bianchi identity, the above equation along with \eqref{eq:tilde-theta-def-1} gives us \eqref{eq:tilde-theta-def-2}.
	
	For the final claim let us assume that $\thetatilde$ is given by \eqref{eq:tilde-theta-def-1} and note that by Theorem~\ref{thm:homogeneity} we can assume without loss of generality that $w=\kappa_s$ for arbitrary $s\in S'$. Then we can use $\fso(V)$-invariance of $\kappa$ to compute
	\begin{align*}
		\bigl(\lambda(u)\cdot\thetatilde\bigr)(v,\kappa_s) + \bigl(\lambda(v)\cdot\thetatilde\bigr)(\kappa_s,u)
			={}& \bigl((\lambda(\lambda(u)v-\lambda(v)u)
				- \comm{\lambda(u)}{\lambda(v)})\cdot\gammahat\bigr)(s,s) \\
				& + 2\gammahat\bigl(s,(\lambda(u)\cdot\betahat)(v,s)
				- (\lambda(v)\cdot\betahat)(u,s)\bigr)\\
				& + 2(\lambda(u)\cdot\gammahat)\bigl(s,\betahat(v,s)\bigr)
				- 2(\lambda(v)\cdot\gammahat)\bigl(s,\betahat(u,s)\bigr).
	\end{align*}
	On the other hand, we recall that $\gammahat(s,s)-\lambda(\kappa_s)\in\fh$, and that $\thetatilde$ is $\fh$-invariant by the first part of the lemma, whence $\lambda(\kappa_s)\cdot\thetatilde=\gammahat(s,s)\cdot\thetatilde$. Using \eqref{eq:tilde-theta-def-1} and a cocycle condition, we find\footnote{Actually, it's necessarily to depolarise \eqref{eq:tilde-theta-def-1} in order to evaluate expressions like $\thetatilde(v,\kappa(s,s'))$ for $v\in V$ $s,s'\in S'$.}
	\begin{align*}
		\bigl(\lambda(\kappa_s)\cdot\thetatilde\bigr)(u,v) ={}& [\gammahat(s,s),\thetatilde(u,w)]
			+ 2\gammahat\bigl(s,\betahat(u,\betahat(v,s))-\betahat(v,\betahat(u,s))\bigr)\\
			& - 2(\lambda(u)\cdot\gammahat)\bigl(s,\betahat(v,s)\bigr)
			+ 2(\lambda(v)\cdot\gammahat)\bigl(s,\betahat(u,s)\bigr).
	\end{align*}
	Putting the two calculations together and then using \eqref{eq:tilde-theta-act-s}, we find
	\begin{gather*}%\label{eq:lambda-bianchi-proof}
		\bigl(\lambda(u)\cdot\thetatilde\bigr)(v,\kappa_s) + \bigl(\lambda(v)\cdot\thetatilde\bigr)(\kappa_s,u) + \bigl(\lambda(\kappa_s)\cdot\thetatilde\bigr)(u,v)\\
		\qquad = \bigl((\lambda(\lambda(u)v-\lambda(v)u) - \comm{\lambda(u)}{\lambda(v)})\cdot\gammahat\bigr)(s,s)
		-\big[\thetatilde(u,w),\gammahat(s,s)\big] \\
		\phantom{\qquad =}{} + 2\gammahat\bigl(s,\betahat\bigl(u,\betahat(v,s)\bigr)-\betahat\bigl(v,\betahat(u,s)\bigr)\bigr)+ 2\gammahat\bigl(s,\bigl(\lambda(u)\cdot\betahat\bigr)(v,s) - \bigl(\lambda(v)\cdot\betahat\bigr)(u,s)\bigr)\\
		\qquad= -\bigl(\bigl(\thetatilde(u,w) + \comm{\lambda(u)}{\lambda(v)} - \lambda(\lambda(u)v-\lambda(v)u)\bigr)\cdot\gammahat\bigr)(s,s).
	\end{gather*}
	The final line vanishes by Lemma~\ref{lemma:theta-lambda-in-h} and $\fh$-invariance of $\gammahat$.
\end{proof}

Thus the problem of integrating an admissible cocycle reduces to the question of whether a~map satisfying the equations \eqref{eq:tilde-theta-def-1} and \eqref{eq:tilde-theta-act-s} exists. Let us first consider equation \eqref{eq:tilde-theta-def-1}. We would like to interpret this as the definition of a map $\thetatilde\colon \Wedge^2V\to \fso(V)$, but \textit{a priori} we only have that
\begin{equation}\label{eq:Theta-def}
	\Theta(v,s,s) = 2\gammahat\bigl(s,\betahat(v,s)\bigr) - \bigl(\lambda(v)\cdot\gammahat\bigr)(s,s)
\end{equation}
defines a map $\Theta\colon V\otimes \Odot^2 S' \to \fso(V)$. So we require  that $\Theta$ factorises as follows
\[%\label{eq:Theta-factor-diagram}
\begin{tikzcd}
	V\otimes \Odot^2 S' \ar[rr,"\Theta"]\ar[rd,"\Id\otimes \kappa"'] & & \fso(V)\\
		& V \otimes V \ar[ru,"\thetatilde"'] &
\end{tikzcd}
\]
and that $\thetatilde$ is skew-symmetric. We define the \emph{Dirac kernel} of $S'$ by $\fD=\ker\kappa|_{\Odot^2S'}$ and note that $\Theta$ factorises as above if and only if it annihilates the Dirac kernel, i.e., $\Theta(V\otimes\fD)=0$. We then find that $\thetatilde$ is automatically skew-symmetric.

\begin{Lemma}\label{lemma:theta-well-def}
	Suppose that $\Theta$ annihilates the Dirac kernel and thus factors as described above. Then equation~\eqref{eq:tilde-theta-def-1} defines a map $\thetatilde\colon\Wedge^2V\to \fso(V)$.
\end{Lemma}

\begin{proof}
	The factorisation of $\Theta$ gives us a map $\thetatilde\colon V\otimes V \to \fso(V)$ satisfying \eqref{eq:tilde-theta-def-1}. To show that $\thetatilde$ is alternating, it is sufficient to show that $\thetatilde(\kappa_s,\kappa_s)=0$ for all $s\in S'$. Using the fact that $\gammahat(s,s)-\lambda(\kappa_s)\in\fh$ for $s\in S'$ and $\fh$-invariance of $\gammahat$, we have $\lambda(\kappa_s)\cdot\gammahat = \gammahat(s,s)\cdot\gammahat$, so
	\begin{align*}
		\thetatilde(\kappa_s,\kappa_s)
		&= 2\gammahat\bigl(s,\betahat(\kappa_s,s)\bigr) - \bigl(\gammahat(s,s)\cdot\gammahat\bigr)(s,s)\\
		& = 2\gammahat\bigl(s,\betahat(\kappa_s,s) + \gammahat(s,s)\cdot s\bigr) - \comm{\gammahat(s,s)}{\gammahat(s,s)}
		= 0,
	\end{align*}
	where we have used a cocycle condition in the last equality.
\end{proof}

\begin{Definition}\label{def:integrable-cocycle}
	Let $[\alpha+\beta+\gamma]$ be an admissible cohomology class and, recalling that it does not depend on the choice of cocycle representative, let $\Theta\colon V\otimes\Odot^2S'\to\fso(V) $ be the map defined by equation~\eqref{eq:Theta-def}. We say that $[\alpha+\beta+\gamma]$ is \emph{integrable} if
	\begin{enumerate}\itemsep=0pt
	\item $\Theta$ annihilates the Dirac kernel $\fD$ of $\Odot^2S'$,
	\item the map $\thetatilde\colon\Wedge^2 V\to \fso(V)$ defined by $\thetatilde\circ\kappa=\Theta$ satisfies	\eqref{eq:tilde-theta-act-s}.
	\end{enumerate}
	An admissible cocycle is \emph{integrable} if its cohomology class is.
\end{Definition}

By construction, we have the following refinement of Proposition~\ref{prop:filtered-def-high-susy}.

\begin{Theorem}[integration of admissible cocycles]\label{thm:admiss-coycle-integr}
	Let $\alpha+\beta+\gamma$ be an admissible, integrable Spencer $(2,2)$-cocycle for a highly supersymmetric graded subalgebra $\fa$ of $\fs$ with $i_*(\alpha+\beta+\gamma)=i^*\bigl(\betahat+\gammahat\bigr)+\partial\lambda$ and let $\thetatilde\colon\Wedge^2V\to\fso(V)$ be the map defined by equation~\eqref{eq:tilde-theta-def-1}. Then the brackets
	\begin{gather*}
			\comm{A}{B} = \underbrace{AB-BA}_\fh,\qquad
			\comm{s}{s} = \underbrace{\kappa_s}_V + \underbrace{\gammahat(s,s) - \lambda(\kappa_s)}_\fh,\\
			\comm{A}{v}= \underbrace{Av}_V + \underbrace{\comm{A}{\lambda(v)} - \lambda(Av)}_\fh,\qquad
			\comm{v}{s} = \underbrace{\betahat(v,s) + \lambda(v)\cdot s}_{S'},\qquad
			\comm{A}{s} = \underbrace{A\cdot s}_{S'},
	\\
		\comm{v}{w} = \underbrace{\lambda(v)w-\lambda(w)v}_V + \underbrace{\thetatilde(v,w) - \lambda(\lambda(v)w-\lambda(w)v) + \comm{\lambda(v)}{\lambda(w)}}_\fh \label{eq:admiss-cocycle-brackets}
	\end{gather*}
	define a geometrically realisable filtered deformation of $\fa$. Two such deformations $\fatilde$, $\fatilde'$ are isomorphic if and only if $\bigl(\betahat'+\gammahat'\bigr)-\bigl(\betahat+\gammahat\bigr)\in\cK^{2,2}(\fa_-)$.
\end{Theorem}

Where $\fa$ is a \emph{maximally} supersymmetric graded subalgebra of the (minimal) Poincar\'e superalgebra $\fs$ in 4, 5, 6 or 11 spacetime dimensions, there is no obstruction to an admissible cocycle being integrable \cite{Beckett2021,deMedeiros2016,deMedeiros2018,Figueroa-OFarrill2017}. We already observed that admissible cohomology classes are simply elements of $\ssH^{2,2}(\fs_-;\fa)^{\fa_0}$ if $\fa$ is highly supersymmetric, so \emph{any} such cohomology class defines a realisable filtered deformation in these cases. The question of whether obstructions exist in the highly supersymmetric (but sub-maximal) case is still open even in these cases.
